\begin{textsample}{NEG Dim 1 – 2018\_c6 – Score -6.00 – 2018\_c6/enade\_2018\_serv\_soc\_f2ecfc474e.txt}  \label{ex:f1_neg_191}
CAPÍTULO 6 ANÁLISE TÉCNICA DA PROVA Este capítulo tem por objetivo apresentar o desempenho dos estudantes concluintes de Serviço Social no Enade/2018.
Para isso, foram calculadas as estatísticas básicas da prova em seu todo ( seção 6.1.1 ), bem como as estatísticas relacionadas aos Componentes de Formação Geral ( seção 6.1.2 ) e de Conhecimento Específico da Área ( seção 6.1.3 ).
Dadas as suas características, foram analisadas, em separado, as questões objetivas ( seção 6.2 ) e as questões discursivas ( seção 6.3 ).
Para as questões objetivas, foram disponibilizados os índices de facilidade e de discriminação Ponto-Bisserial, também em separado, para os Componentes de Formação Geral ( seção 6.2.1 ) e de Conhecimento Específico ( seção 6.2.2 ).
De \textbf{cada} componente, uma das questões foi escolhida para exemplificar a análise gráfica, relacionando -se as alternativas escolhidas pelos estudantes ( inclusive o gabarito ) com o número de acertos no componente.
No Anexo I, apresenta -se a íntegra da análise gráfica para todas as questões objetivas.
Para \textbf{cada} uma das questões discursivas, os conteúdos dos tipos mais comuns de respostas dos estudantes são apresentados e comparados com o padrão de respostas esperado ( ver Anexo VIII com o padrão de respostas ).
Tomando -se como base as duas questões discursivas do Componente de Formação Geral, na seção 6.3.1.6 são apresentados comentários sobre a correção das respostas em relação à Língua Portuguesa.
Nas tabelas, constam as seguintes estatísticas das notas20 : média do desempenho na prova, erro padrão da média, desvio padrão, nota mínima, mediana e nota máxima para \textbf{cada} um de seus componentes.
Tais estatísticas contemplam o total de estudantes concluintes da Área de Serviço Social inscritos e presentes à prova do Enade/2018, tendo em vista agregações, ou por Grandes Regiões ou o país como um todo, ou por Categoria Administrativa, Organização Acadêmica da IES e Modalidade de Ensino.
Em relação aos gráficos de distribuição de notas, o intervalo considerado foi de 10 unidades, aberto à esquerda e fechado à direita, com exceção do primeiro intervalo, [ 0 ; 10 ], fechado em ambos os extremos.
Para os gráficos de distribuição das notas das questões discursivas, foram consideradas mais duas categorias : questão em branco21 e nota zero. 20 Uma definição dessas estatísticas pode ser encontrada no Glossário. 21 Nesse grupo \textbf{estão} incluídas também as respostas classificadas como nulas ou desconsideradas. 6.1 ESTATÍSTICAS BÁSICAS DA PROVA Esta seção apresenta estatísticas selecionadas e histogramas da nota geral ( 6.1.1 ) e de \textbf{cada} componente : Formação Geral ( 6.1.2 ) e Conhecimento Específico ( 6.1.3 ).
São também apresentadas estatísticas selecionadas de subpopulações, caracterizadas por Grande Região, Categoria Administrativa, Organização Acadêmica e Modalidade de Ensino. 6.1.1 Estatísticas Básicas Gerais Na Tabela 6.1, são apresentadas as Estatísticas Básicas da prova, por Grande Região, dos estudantes concluintes de Serviço Social.
A Média das notas da prova, como um todo ( nas seções seguintes serão analisados os Componentes de Formação Geral e de Conhecimento Específico ) foi 40,0, sendo que os alunos da região Sul obtiveram a Média mais baixa ( 36,1 ), e os da região Nordeste obtiveram a Média mais alta ( 44,6 ).
As demais médias foram : 42,0, na região Norte, 40,4, na região Sudeste e 38,3, na região Centro-Oeste.
O Desvio padrão para o Brasil, como um todo, foi 15,8, sendo o maior Desvio padrão encontrado na região Nordeste ( 16,4 ), e o menor, na região Sul ( 13,7 ), indicando uma dispersão um pouco menor das notas dessa última região.
A região que obteve a maior nota Máxima foi a Nordeste ( 91,2 ), ao passo que as regiões que atingiram a menor nota Máxima foram a Norte e a Sudeste ( 90,3 ).
A Mediana do Brasil, como um todo, foi 37,9, sendo a maior Mediana obtida na região Nordeste ( 43,5 ), e a menor, obtida na Sul ( 34,3 ).
A nota Mínima do Brasil, zero foi obtida em quase todas as regiões, exceto na região Norte ( 3,7 ).
Considerando -se as notas segundo Grande Região, observa -se que existe diferença estatisticamente significativa ao nível de 95 \% entre todas as regiões22. 22 Todas as comparações deste capítulo utilizam os intervalos de 95 \%.
Os erros-padrão da média que possibilitam os testes \textbf{estão} disponíveis nas tabelas.
O comportamento das notas dos estudantes de todo o Brasil pode ser observado no Gráfico 6.1, que apresenta um histograma com a distribuição das mesmas : uma distribuição unimodal com a moda no intervalo ( 30 ; 40 ], 25,2 \%, seguido de perto pelo intervalo ( 20 ; 30 ], 21,8 \%.
Na Tabela 6.2, são apresentadas informações referentes à Média da nota final, desagregadas por Categoria Administrativa, por Organização Acadêmica e por Modalidade de Ensino.
A Média das notas da prova, como um todo, foi 40,0.
Em relação à Categoria Administrativa, os estudantes das IES Públicas obtiveram Média bem mais alta ( 53,9 ) que a Média nacional, e os das IES Privadas obtiveram Média mais baixa ( 37,3 ).
Observa -se que existe diferença estatisticamente significativa entre as médias das notas das IES Públicas e as das Privadas.
A diferença entre as médias das regiões Nordeste e Sul ( 8,5 ), a maior e a menor Média, é inferior à diferença entre a das IES Públicas e a das Privadas ( 16,6 ), caracterizando -se uma maior diversidade administrativa do que regional.
O Desvio padrão para as IES Públicas ( 16,4 ) foi superior ao do Brasil, como um todo ( 15,8 ), indicando uma dispersão um pouco maior das notas nesta Categoria Administrativa.
No tocante à Organização Acadêmica, os CEFET/IFET obtiveram Média mais alta que a nacional ( 61,0 ).
A Média das Faculdades e dos Centros Universitários foram mais alta que a nacional, também, no entanto bem mais próxima a ela ( 40,8 e 40,1, respectivamente ).
A Média das Universidades ( 39,7 ) foi inferior à nacional.
Constata -se que existe diferença estatisticamente significativa ao nível de 95 \% entre a Média mais alta, 60,1 dos CEFET/IFET, e a dos demais tipos de Organização Acadêmica.
A Média da Modalidade Educação a Distância ( 33,6 ) foi inferior à Média da Educação Presencial ( 46,2 ), e a diferença entre elas é estatisticamente significativa, ao nível de 95 \%. 6.1.2 Estatísticas Básicas no Componente de Formação Geral Na Tabela 6.3, são apresentadas as Estatísticas Básicas em relação ao componente da prova que avalia a Formação Geral dos estudantes concluintes.
Os alunos de todo o Brasil obtiveram desempenho médio de 40,1.
Quanto à variabilidade, o Desvio padrão das notas dos estudantes do Brasil, como um todo, foi 16,6.
A maior Média foi obtida na região Nordeste ( 42,0 ), e a menor, na região Centro-Oeste ( 38,3 ).
As demais médias foram : 41,2, na região Norte, 40,4, na região Sudeste e 38,7, na região Sul.
Já o maior Desvio padrão foi obtido na região Centro-Oeste ( 17,1 ), e o menor, na região Sul ( 15,7 ).
Os demais desvios padrões foram : 16,7, na região Norte, 16,8 na região Nordeste e 16,9 na região Sudeste.
A maior nota no Componente de Formação Geral da prova do Enade/2018 foi 93,8, obtida por, pelo menos, um aluno na região Sudeste.
A menor nota Máxima foi obtida na região Norte ( 90,2 ).
A Mediana do Brasil, como um todo, foi 39,2, sendo a menor Mediana encontrada na região Centro-Oeste ( 37,5 ), e a maior, na região Nordeste ( 41,3 ).
A nota Mínima nessa parte foi zero em todas as regiões.
Considerando -se as notas, segundo a Grande Região, observa -se que não existe diferença estatisticamente significativa entre a Média mais alta ( região Nordeste ) e a da região Norte, mas são estatisticamente significativas as diferenças entre a Média da região Nordeste e das demais regiões.
No Gráfico 6.2, pode -se verificar a avaliação do desempenho dos estudantes do Componente de Formação Geral, a partir do histograma da distribuição das notas correspondentes.
A distribuição é unimodal, com moda em ( 30 ; 40 ], o mesmo intervalo modal da distribuição de notas da prova, como um todo ( Gráfico 6.1 ).
Nota -se, ainda, que, no Gráfico 6.2, as notas apresentam maior dispersão do que as no Gráfico 6.1 ( distribuição das notas da prova ), confirmada pela comparação dos desvios padrão : 15,8 para a nota da prova, como um todo, e 16,6 para o Componente de Formação Geral.
Na Tabela 6.4, são apresentadas as informações referentes ao desempenho dos concluintes do Componente de Formação Geral, em diferentes agregações : Categoria Administrativa, Organização Acadêmica e Modalidade de Ensino.
Observa -se que existe diferença estatisticamente significativa entre as médias dos tipos de Categoria Administrativa : a maior Média obtida por estudantes de IES Públicas ( 49,1 ) e a menor, pelos de IES Privadas ( 38,3 ).
Nota -se que, também, há diferença estatisticamente significativa entre a maior das médias dos quatro tipos de Organização Acadêmica e das demais, o valor maior para os CEFET/IFET ( 53,9 ) e a menor para as Faculdades ( 39,3 ).
No que se refere à Modalidade de Ensino, a Educação Presencial apresentou maior Média do que a Educação a Distância ( médias 43,4 e 36,7, respectivamente ).
Tais médias também apresentam diferença estatisticamente significativa. 6.1.3 - Estatísticas Básicas do Componente de Conhecimento Específico Na Tabela 6.5, são apresentadas as Estatísticas Básicas referentes ao Componente de Conhecimento Específico da área de Serviço Social.
A Média do desempenho dos alunos do Brasil, como um todo, foi 40,0, a mesma da prova como um todo.
A maior Média foi obtida na região Nordeste ( 45,5 ), e a menor, na região Sul ( 35,2 ).
As demais médias foram : 42,3, na região Norte, 40,4, na região Sudeste e 38,3, na região Centro-Oeste.
Quanto à variabilidade das notas, o Desvio padrão do Brasil, como um todo, foi 17,3, sendo o maior Desvio padrão observado na região Nordeste ( 18,0 ), e o menor, na região Sul ( 14,8 ).
Os demais desvios foram : 16,4, na região Norte, 17,8, na região Sudeste e 17,1, na região Centro-Oeste.
A Mediana das notas dos estudantes de todo o Brasil foi 37,7, o mesmo valor da Mediana do Sudeste.
A maior Mediana ocorreu na região Nordeste ( 44,1 ), e a menor, na região Sul ( 33,0 ).
As demais medianas foram : 40,7, na região Norte, 37,7, na região Sudeste e 34,9, na região Centro-Oeste.
A nota Máxima do Brasil, como um todo, foi 97,5, sendo obtida por, pelo menos, um aluno na região Sudeste.
As demais notas máximas foram : 90,7 na região Norte, 92,9, na região Nordeste, 91,4, na região Sul e 90,9, na região Centro-Oeste.
A nota Mínima foi zero em todas as regiões.
Observa -se que existe diferença estatisticamente significativa entre as Médias das notas do Componente de Conhecimento Específico de todas as regiões.
Assim como nos Gráficos 6.1 e 6.2, no Gráfico 6.3, apresentado a seguir, é possível ser feita uma avaliação do desempenho de concluintes de Serviço Social em relação ao Componente de Conhecimento Específico, com um histograma da distribuição das notas correspondentes.
Essa também é uma distribuição unimodal, e o grupo modal é o ( 30 ; 40 ], o mesmo grupo modal da prova, como um todo, e da Formação Geral.
Na Tabela 6.6, é apresentada uma comparação dos resultados em relação à Categoria Administrativa, à Organização Acadêmica e à Modalidade de Ensino, agora levando -se em conta o desempenho de alunos do Componente de Conhecimento Específico da prova.
No que se refere à Organização Acadêmica, a maior Média foi a dos CEFET/IFET ( 63,3 ), vindo a seguir a das Faculdades ( 41,3 ), a dos Centros Universitários ( 40,0 ) e, por fim, a das Universidades ( 39,6 ).
O maior Desvio padrão, e acima do valor para o Brasil, como um todo, foi o das Universidades ( 17,6 ).
As Universidades obtiveram a maior nota Máxima ( 97,5 ).
Os Centros Universitários obtiveram nota Máxima 94,9, os CEFET/IFET, 86,7 e as Faculdades obtiveram nota Máxima 88,4.
As medianas foram : 36,6 nas Universidades, a menor delas ; 38,6 nos Centros Universitários ; 40,1 nas Faculdades ; e 65,5 nos CEFET/IFET, a maior.
A nota Mínima foi zero para quase todas as Organizações Acadêmicas, exceto para os CEFET/IFET ( 22,0 ).
Observa -se que existe diferença estatisticamente significativa ao nível de 95 \% no Componente de Conhecimento Específico entre as notas dos CEFET/IFET e das demais Organizações Acadêmicas.
Quanto à Categoria Administrativa, observa -se um comportamento semelhante àquele da parte de Formação Geral e à prova, como um todo, ou seja, existe diferença 6.2 ANÁLISE DAS QUESTÕES OBJETIVAS Esta seção apresenta estatísticas selecionadas e histogramas das Questões Objetivas de Formação Geral ( 6.2.1 ) e de Conhecimento Específico ( 6.2.2 ).
São também apresentadas e comparadas as médias das subpopulações caracterizadas por Grande Região. 6.2.1 Componente de Formação Geral Na Tabela 6.7, são apresentadas as Estatísticas Básicas relativas às oito questões objetivas do componente da prova que abrange a Formação Geral dos estudantes.
A Média do Brasil foi 43,6.
A menor Média foi encontrada na região Sul ( 41,7 ), e a maior, na região Nordeste ( 45,6 ).
As demais médias foram : 44,2, na região Norte, 43,9, na região Sudeste e 43,5, na região Centro-Oeste.
O Desvio padrão do Brasil foi 21,5, sendo o maior Desvio padrão encontrado nas regiões Sudeste e Centro-Oeste ( 21,9 ), e o menor, na região Sul ( 20,7 ).
Os demais desvios foram : 21,2, na região Norte e 21,6, na região Nordeste.
As medianas do Brasil, como um todo, e de quatro das cinco regiões foram 37,5 - a única região com Mediana diferente foi a Nordeste, com 50,0.
A nota Máxima 100,0, foi alcançada em todas as Grandes Regiões.
As notas mínimas ( 0,0 ) foram iguais para todas as regiões. estatisticamente significativa entre as médias das IES Públicas ( 55,5 ) e as das IES Privadas ( 36,9 ).
Nesse caso, também, a maior Média foi obtida por alunos de IES Públicas de ensino.
Quanto ao comportamento das médias no que toca à Modalidade de Ensino, há diferença estatisticamente significativa entre a Média da Educação Presencial ( 47,1 ) e a da Educação a Distância ( 32,5 ). 6.2 ANÁLISE DAS QUESTÕES OBJETIVAS Esta seção apresenta estatísticas selecionadas e histogramas das Questões Objetivas de Formação Geral ( 6.2.1 ) e de Conhecimento Específico ( 6.2.2 ).
São também apresentadas e comparadas as médias das subpopulações caracterizadas por Grande Região. 6.2.1 Componente de Formação Geral Na Tabela 6.7, são apresentadas as Estatísticas Básicas relativas às oito questões objetivas do componente da prova que abrange a Formação Geral dos estudantes.
A Média do Brasil foi 43,6.
A menor Média foi encontrada na região Sul ( 41,7 ), e a maior, na região Nordeste ( 45,6 ).
As demais médias foram : 44,2, na região Norte, 43,9, na região Sudeste e 43,5, na região Centro-Oeste.
O Desvio padrão do Brasil foi 21,5, sendo o maior Desvio padrão encontrado nas regiões Sudeste e Centro-Oeste ( 21,9 ), e o menor, na região Sul ( 20,7 ).
Os demais desvios foram : 21,2, na região Norte e 21,6, na região Nordeste.
As medianas do Brasil, como um todo, e de quatro das cinco regiões foram 37,5 - a única região com Mediana diferente foi a Nordeste, com 50,0.
A nota Máxima 100,0, foi alcançada em todas as Grandes Regiões.
As notas mínimas ( 0,0 ) foram iguais para todas as regiões.
Na Tabela 6.8, são apresentados o Índice de Facilidade e o Índice de Discriminação ( Ponto-Bisserial ) para \textbf{cada} uma das questões objetivas do Componente de Formação Geral.
Quanto ao Índice de Facilidade, foram usadas as seguintes cores para diferenciar o nível de dificuldade da questão :  Azul para as questões classificadas com índice Muito fácil ( \&gt ; =0,86 ), verde para as questões classificadas com índice Fácil ( 0,61 a 0,85 ), amarelo para as questões classificadas com Médio ( 0,41 a 0,60 ), vermelho para as questões classificadas com Difícil ( 0,16 a 0,40 ) e roxo para as questões classificadas com Muito difícil ( \&lt ; =0,15 ).
Já quanto ao Índice de Discriminação, foram usadas as seguintes cores para qualificar a questão :  As questões classificadas com índice Fraco receberam a cor vermelha ( \&lt ; =0,19 ), as classificadas com Médio receberam a cor amarela ( 0,20 a 0,29 ), as classificadas com Bom receberam a cor verde ( 0,30 a 0,39 ) e as classificadas com Muito bom ( \&gt ; =0,40 ) receberam a cor azul.
As questões objetivas do Componente de Formação Geral, segundo o Índice de Facilidade, foram assim avaliadas : das oito questões, nenhuma teve o Índice de Facilidade classificado como Muito fácil, e duas questões foram consideradas com índice de dificuldade Fácil, situando -se no intervalo entre 0,61 e 0,85 do Índice de Facilidade, ou seja, houve entre 61,0 \% e 85,0 \% de acertos.
Três questões foram consideradas com índice de dificuldade Médio, situando -se no intervalo entre 0,41 e 0,60 do Índice de Facilidade, ou seja, houve entre 41,0 \% e 60,0 \% de acertos, enquanto outras três questões foram classificadas na categoria Difícil, situando -se no intervalo entre 0,16 e 0,40.
Nenhuma questão apresentou menos de 15 \% de acertos, razão pela qual não houve questão classificada como Muito difícil.
O Índice de Facilidade variou de 0,22 a 0,63, e o de Discriminação, de 0,42 a 0,49.
Na Tabela 6.8, são apresentados o Índice de Facilidade e o Índice de Discriminação ( Ponto-Bisserial ) para \textbf{cada} uma das questões objetivas do Componente de Formação Geral.
Quanto ao Índice de Facilidade, foram usadas as seguintes cores para diferenciar o nível de dificuldade da questão :  Azul para as questões classificadas com índice Muito fácil ( \&gt ; =0,86 ), verde para as questões classificadas com índice Fácil ( 0,61 a 0,85 ), amarelo para as questões classificadas com Médio ( 0,41 a 0,60 ), vermelho para as questões classificadas com Difícil ( 0,16 a 0,40 ) e roxo para as questões classificadas com Muito difícil ( \&lt ; =0,15 ).
Já quanto ao Índice de Discriminação, foram usadas as seguintes cores para qualificar a questão :  As questões classificadas com índice Fraco receberam a cor vermelha ( \&lt ; =0,19 ), as classificadas com Médio receberam a cor amarela ( 0,20 a 0,29 ), as classificadas com Bom receberam a cor verde ( 0,30 a 0,39 ) e as classificadas com Muito bom ( \&gt ; =0,40 ) receberam a cor azul.
As questões objetivas do Componente de Formação Geral, segundo o Índice de Facilidade, foram assim avaliadas : das oito questões, nenhuma teve o Índice de Facilidade classificado como Muito fácil, e duas questões foram consideradas com índice de dificuldade Fácil, situando -se no intervalo entre 0,61 e 0,85 do Índice de Facilidade, ou seja, houve entre 61,0 \% e 85,0 \% de acertos.
Três questões foram consideradas com índice de dificuldade Médio, situando -se no intervalo entre 0,41 e 0,60 do Índice de Facilidade, ou seja, houve entre 41,0 \% e 60,0 \% de acertos, enquanto outras três questões foram classificadas na categoria Difícil, situando -se no intervalo entre 0,16 e 0,40.
Nenhuma questão apresentou menos de 15 \% de acertos, razão pela qual não houve questão classificada como Muito difícil.
O Índice de Facilidade variou de 0,22 a 0,63, e o de Discriminação, de 0,42 a 0,49.
Na Tabela 6.9, é apresentada a distribuição das questões, levando -se em conta, simultaneamente, a classificação dos dois índices.
Todas as questões obtiveram Índice de Discriminação Muito bom e figuraram entre três níveis de dificuldade Fácil, Médio e Difícil : duas classificadas na categoria Fácil ( questões 2 e 8 ) do Índice de Facilidade, três na categoria Médio ( questões 3, 6 e 7 ), e outras três, na categoria Difícil ( questões 1, 4 e 5 ).
Em particular, a questão 5, foi uma das que apresentou o maior poder discriminatório, com índice 0,49, e foi considerada Difícil em termos de facilidade, com uma proporção de 0,40 acertos.
O máximo de acertos foi alcançado pela questão 8 com um Índice de Facilidade de 0,63.
As questões 1 e 4 foram consideradas Difícil, com índices de Facilidade 0,27 e 0,22, respectivamente.
Já os índices de Discriminação foram Muito bom para as duas questões ( 0,43 para \textbf{cada} uma delas ).
No Gráfico 6.4, para exemplificar, analisa -se o comportamento da questão de número 4 de Formação Geral.
Trata -se de uma questão considerada Médio, em relação à facilidade e a que obteve o segundo maior Índice de Discriminação dessa parte da prova.
Nesse gráfico, \textbf{cada} uma das cinco curvas representa o percentual de respostas em determinada alternativa da questão, em função do número de acertos dos estudantes nessa Na Tabela 6.9, é apresentada a distribuição das questões, levando -se em conta, simultaneamente, a classificação dos dois índices.
Todas as questões obtiveram Índice de Discriminação Muito bom e figuraram entre três níveis de dificuldade Fácil, Médio e Difícil : duas classificadas na categoria Fácil ( questões 2 e 8 ) do Índice de Facilidade, três na categoria Médio ( questões 3, 6 e 7 ), e outras três, na categoria Difícil ( questões 1, 4 e 5 ).
Em particular, a questão 5, foi uma das que apresentou o maior poder discriminatório, com índice 0,49, e foi considerada Difícil em termos de facilidade, com uma proporção de 0,40 acertos.
O máximo de acertos foi alcançado pela questão 8 com um Índice de Facilidade de 0,63.
As questões 1 e 4 foram consideradas Difícil, com índices de Facilidade 0,27 e 0,22, respectivamente.
Já os índices de Discriminação foram Muito bom para as duas questões ( 0,43 para \textbf{cada} uma delas ).
No Gráfico 6.4, para exemplificar, analisa -se o comportamento da questão de número 4 de Formação Geral.
Trata -se de uma questão considerada Médio, em relação à facilidade e a que obteve o segundo maior Índice de Discriminação dessa parte da prova.
Nesse gráfico, \textbf{cada} uma das cinco curvas representa o percentual de respostas em determinada alternativa da questão, em função do número de acertos dos estudantes nessa Na Tabela 6.9, é apresentada a distribuição das questões, levando -se em conta, simultaneamente, a classificação dos dois índices.
Todas as questões obtiveram Índice de Discriminação Muito bom e figuraram entre três níveis de dificuldade Fácil, Médio e Difícil : duas classificadas na categoria Fácil ( questões 2 e 8 ) do Índice de Facilidade, três na categoria Médio ( questões 3, 6 e 7 ), e outras três, na categoria Difícil ( questões 1, 4 e 5 ).
Em particular, a questão 5, foi uma das que apresentou o maior poder discriminatório, com índice 0,49, e foi considerada Difícil em termos de facilidade, com uma proporção de 0,40 acertos.
O máximo de acertos foi alcançado pela questão 8 com um Índice de Facilidade de 0,63.
As questões 1 e 4 foram consideradas Difícil, com índices de Facilidade 0,27 e 0,22, respectivamente.
Já os índices de Discriminação foram Muito bom para as duas questões ( 0,43 para \textbf{cada} uma delas ).
No Gráfico 6.4, para exemplificar, analisa -se o comportamento da questão de número 4 de Formação Geral.
Trata -se de uma questão considerada Médio, em relação à facilidade e a que obteve o segundo maior Índice de Discriminação dessa parte da prova.
Nesse gráfico, \textbf{cada} uma das cinco curvas representa o percentual de respostas em determinada alternativa da questão, em função do número de acertos dos estudantes nessa parte da prova ( Formação Geral/Múltipla Escolha ), antes de possíveis eliminações pelo critério Ponto-Bisserial.
A curva em verde corresponde à alternativa B, a alternativa correta para essa questão.
Observa -se que entre os estudantes com menor número de acertos, nessa parte do exame, a situação mais frequente foi a escolha da alternativa C ( em preto ), incorreta.
Por exemplo, entre os estudantes que acertaram três questões, 36,1 \%, escolheram a alternativa C ( em preto ), 20,7 \%, a alternativa E ( em vermelho ), 20,0 \%, escolheram a alternativa D ( em laranja ), 13,0 \%, a alternativa B ( em verde ) e 10,0 \%, a alternativa A ( azul ).
Entre os que acertaram três respostas das questões de múltipla escolha de Formação Geral, 0,1 \%, deixou a questão em branco e 0,2 \%, marcou mais de uma alternativa, invalidando a questão.
À medida que o número de acertos aumenta, indicando desempenho melhor nessa parte da prova, aumenta concomitantemente a proporção de estudantes que selecionaram a alternativa correta B, atingindo 100 \%, para os estudantes com oito acertos.
Essa análise permite verificar como a questão discriminou os grupos de desempenho, justificando -se o alto índice obtido na questão.
Cumpre notar que não é possível inferir deste gráfico nem o índice de facilidade ( que seria uma média da proporção ponderada pela quantidade de alunos com \textbf{cada} uma das notas ), nem o índice de discriminação Ponto-Bisserial, por razão equivalente.
No caso extremo, no qual a grande concentração dos acertos dos alunos fosse abaixo de quatro, o índice de facilidade seria obrigatoriamente abaixo de 20 \% ( neste exemplo ).
Caso a concentração fosse em seis acertos ou mais, o índice seria obrigatoriamente acima de 50 \%.
Os gráficos relativos às demais questões de Formação Geral constam do Anexo I. 6.2.2 Componente de Conhecimento Específico Na Tabela 6.10, são apresentadas as Estatísticas Básicas em relação às questões objetivas do Componente de Conhecimento Específico da prova de Serviço Social, por Grande Região.
Nessa parte da prova, nenhuma questão foi anulada pela Comissão Assessora da Área.
Além disso, como será discutido ainda nesta seção, nenhuma questão objetiva do Componente de Conhecimentos Específicos não foi usada no cômputo das notas por ter sido descartada pelo critério Ponto-Bisserial.
Assim, as notas foram calculadas com base nas 27 das 27 questões objetivas de Conhecimentos Específicos.
A Média do Brasil deste componente foi 43,2.
A menor Média foi observada na região Sul ( 38,1 ), e a maior, na região Nordeste ( 49,0 ).
O Desvio padrão de todo o Brasil foi 18,4, sendo o menor Desvio padrão encontrado na região Sul ( 16,0 ), e o maior, na região Nordeste ( 19,1 ).
A Mediana de todo o Brasil foi 40,7, o mesmo valor da Mediana encontrada na região Sudeste.
Nas regiões Sul e Centro-Oeste, a Mediana foi menor : 37,0, em \textbf{cada} uma dessas regiões.
A nota Máxima da prova foi 100,0, obtida nas questões objetivas do Componente de Conhecimento Específico, por, pelo menos, um aluno da região Sudeste.
A nota Mínima foi zero em todas as regiões.
Na Tabela 6.11, são apresentados os Índices de Facilidade e Discriminação ( Ponto-Bisserial ) das questões objetivas do Componente de Conhecimento Específico, para os estudantes de Serviço Social.
Para facilitar a diferenciação das questões, serão usadas as mesmas cores da Tabela 6.8 para as diferentes classificações dos Índices de Facilidade e de Discriminação.
Observa -se que nenhuma questão foi anulada pela CAA, e que as análises dizem respeito às 27 questões válidas.
A partir do Índice de Facilidade obtido, pode -se concluir que quase 45 \% das questões objetivas da prova foram considerados, pelo menos, Difícil : das 27 questões válidas, 12 foram classificadas como Difícil.
Duas questões foram classificadas como Fácil, e outras 13 consideradas como Médio.
Já quanto aos índices de discriminação das questões objetivas do Componente de Conhecimento Específico da prova, obtém -se, como resultado, a seguinte classificação : 24 das 27 questões foram consideradas boas, enquanto 14 delas tiveram Índice de Discriminação Muito bom.
Assim, para 90 \% das questões – 24 em 27 – os Índices de Discriminação foram Bom ou Muito bom.
Dentre as demais, três delas foram classificadas como Médio, por conseguinte, a quantidade de questões nos patamares mais baixos de discriminação.
Constata -se, assim, que a prova – no que se refere ao Componente de Conhecimento Específico – possuía boa capacidade de discriminar dentre aqueles que dominam ou não o conteúdo.
O Índice de Facilidade variou de 0,16 a 0,85, e o de Discriminação, de 0,21 a 0,51.
Na Tabela 6.12, é apresentada a distribuição das questões levando -se em conta, simultaneamente, a classificação dos dois índices.
Entre as questões que alcançaram o Índice de Discriminação Muito bom, cinco foram classificadas como Difícil quanto ao Índice de Facilidade, oito como Médio e uma como Fácil.
A questão de número 33 foi uma das mais difíceis, pois seu Índice de Facilidade foi 0,40, apenas 40 \% de acertos, e seu Índice de Discriminação foi 0,49.
Por outro lado, a questão de número 26 foi uma das mais fáceis, pois seu Índice de Facilidade foi 0,72, ou seja, 72 \% de acertos com 0,42 de Índice de Discriminação.
Dez questões foram classificadas quanto ao Índice de Discriminação como Bom, e quanto ao Índice de Facilidade, foram cinco em Difícil, quatro em Médio e uma em Fácil.
A questão de número 25 foi a mais fácil das questões com Índice de Facilidade 0,85 e 0,37 de Índice de Discriminação.
Entre as três questões com Índice de Discriminação Médio, duas foram classificadas quanto ao Índice de Facilidade em Difícil e uma em Médio.
A questão de número 20 foi a mais difícil das questões, pois seu Índice de Facilidade foi 0,16, que significa que houve apenas 16 \% de acertos e 0,21 de Índice de Discriminação.
Não houve questões eliminadas pelo critério Ponto-Bisserial.
A título de exemplo das análises do comportamento das questões objetivas, no Gráfico 6.5, analisa -se a questão 14 do Componente de Conhecimento Específico.
Essa questão foi considerada, pelas respostas dos estudantes avaliados na prova, como uma questão com nível de facilidade Médio, 0,51, ou seja, 51 \% dos estudantes assinalaram, acertadamente, a opção C, correspondente ao gabarito.
Como já comentado, seu Índice de Discriminação foi igual a 0,51, classificado como Muito bom, o maior valor de discriminação.
Nesse gráfico, \textbf{cada} uma das cinco curvas representa o percentual de respostas em determinada alternativa da questão 14, em função do número de acertos dos estudantes nessa parte da prova, antes de possíveis eliminações de questões pelo critério Ponto-Bisserial.
A alternativa correta C, representada no gráfico pela curva em preto, foi escolhida em maiores proporções pelos alunos com desempenho melhor nessa parte da prova.
Já as alternativas incorretas, também denominadas distratores, foram selecionadas, principalmente, por aqueles com notas mais baixas.
Observa -se que a soma não é 100 \%, por causa das questões não respondidas ou com mais de uma opção marcada.
Aqueles com nota zero, na sua quase totalidade, escolheram uma alternativa incorreta, deixaram essa questão em branco ou marcaram mais de uma alternativa, comportamento considerado inválido.
A proporção de alunos que selecionaram a resposta correta C aumenta gradativamente, chegando a atingir 100 \% para 25 acertos ou mais, enquanto a proporção dos que escolheram alternativas incorretas decai, a partir de dois ou três acertos, em função do número de acertos nessa parte da prova.
Os gráficos relativos às demais questões do Conhecimento Específico constam do Anexo I. 6.3 ANÁLISE DAS QUESTÕES DISCURSIVAS Esta seção apresenta estatísticas selecionadas e histogramas das Questões Discursivas de Formação Geral ( 6.3.1 ) e Conhecimento Específico ( 6.3.2 ).
São também apresentadas e comparadas as médias de subpopulações, caracterizadas por Grande Região. 6.3.1 Componente de Formação Geral As análises dos resultados de desempenho dos estudantes de Serviço Social nas duas questões discursivas relativas à Formação Geral encontram -se na Tabela 6.13 e no Gráfico 6.6.
Na Tabela 6.13, observa -se que a nota Média nesse conjunto de questões foi abaixo da obtida nas objetivas.
Os estudantes de todo Brasil obtiveram, em Formação Geral, Média 43,6 nas questões objetivas e 34,8, nas questões discursivas.
Pode -se notar, também, que o Desvio padrão nesse conjunto de questões foi menor do que o obtido nas objetivas : 21,5, nas questões objetivas e 19,1, nas questões discursivas.
A maior Média foi obtida na região Norte ( 36,8 ), e a menor, na região Centro-Oeste ( 30,4 ).
A Mediana de todo o Brasil, neste componente, foi 37,0.
Nas regiões Sul e Centro-Oeste, a Mediana foi menor ( 36,0 e 32,0, respectivamente ), e nas demais regiões foi maior que o valor para o Brasil : Norte ( 40,0 ), Nordeste ( 39,5 ) e Sudeste ( 37,5 ).
A nota Máxima ( 92,5 ) foi obtida na região Sudeste, sendo 84,5 a Máxima nas regiões Norte e Centro-Oeste ; 87,5 na Nordeste e 90,5 na Sul.
A nota Mínima ( 0,0 ) foi a mesma em todas as regiões do Brasil.
No Gráfico 6.6, \textbf{está} representada a distribuição das notas nas questões discursivas do Componente de Formação Geral.
A moda dessa distribuição ocorre no intervalo ( 40 ; 50 ], 21,3 \%, seguida pelo intervalo ( 30 ; 40 ], 20,9 \%.
Destaca -se também, como um máximo local, o intervalo [ 0 ; 10 ], com distribuição de 14,3 \% do total de notas, sendo que no intervalo [ 0 ; 10 ], inclui -se, além da nota zero, a frequência de alunos que deixaram esse tipo de questão em branco.
Na sequência, os resultados verificados para \textbf{cada} uma das questões discursivas de Formação Geral serão apresentados, estabelecendo -se relações com os conteúdos abordados em \textbf{cada} uma delas.
Os comentários da Banca de docentes corretores, a respeito do observado na correção das respostas dos estudantes, suas impressões e conclusões serão apresentados junto com a análise de \textbf{cada} questão.
Cumpre esclarecer que, tendo -se em vista que as questões discursivas de Formação Geral são padronizadas, ou seja, constam de todas as provas, os comentários da Banca são os mesmos para todas as carreiras acadêmicas, sendo direcionados a todos os estudantes que participaram do Enade/2018.
A seguir, serão analisados os desempenhos da Área de Serviço Social nas duas questões discursivas de Formação Geral do Enade/2018, comparando -se os resultados obtidos com comentários para \textbf{cada} questão. 6.3.1.1 Análise de Conteúdo da Questão Discursiva 1 do Componente de Formação Geral Os dados de Serviço Social, obtidos a partir das respostas à questão 1, encontram -se na Tabela 6.14 e no Gráfico 6.7.
Nessa questão – a de melhor desempenho entre as duas de Formação Geral – os alunos de todo o Brasil tiveram Média 31,2.
A maior Média para a questão 1 foi obtida na região Nordeste ( 33,3 ), e a menor, na região Centro-Oeste ( 26,4 ).
Quanto à variabilidade das notas, o Desvio padrão de todo o Brasil foi 23,0.
O menor Desvio padrão foi obtido na região Sul ( 21,5 ), e o maior Desvio padrão foi obtido na região Nordeste ( 24,0 ).
A nota Máxima, a Mediana e a nota Mínima foram as mesmas para todas as regiões, 100,0, 30,0 e zero, respectivamente.
No Gráfico 6.7, mostra -se a distribuição das notas na questão discursiva 1 do Componente de Formação Geral.
Observa -se que a maior frequência corresponde ao intervalo ( 40 ; 50 ], 19,5 \%, sendo seguido de perto pelo intervalo anterior.
Cerca de 11,5 \% dos participantes deixaram a resposta a essa questão em branco, e quase 9 \% obtiveram zero. 6.3.1.2 Comentários sobre a correção de Conteúdo das respostas à Questão Discursiva 1 A Questão 1 exigia do aluno um posicionamento sobre conteúdo humanista e crítico.
O comando da questão solicitava a redação de um texto que formulasse e articulasse argumentos consistentes envolvendo temas previstos na Portaria nº 444, de 30 de maio de 2018, que dispõe sobre o componente de Formação geral do Enade 2018, tais como : solidariedade/violência, tolerância/intolerância, inclusão/exclusão, sexualidade, relações de gênero e relações étnico-raciais ( art. 7º, item VIII ), bem como ética, democracia e cidadania ( art. 7º, item I ).
O enunciado da questão buscava o estabelecimento da relação de três dimensões – os Direitos Humanos, a Democracia e as Fake News.
Para isso, recorreu a quatro suportes : a reprodução de uma foto e três fragmentos adaptados de textos que foram disponibilizados em páginas da Internet ( da Anistia Internacional, do Jornal El País-Brasil e do Labic-Laboratório de estudos sobre Imagem e Cibercultura ).
Os fragmentos traziam dados sobre o aumento do número de assassinatos de grupos minoritários, as origens e o trabalho da vereadora Marielle Franco, sua luta pelos Direitos Humanos, seu assassinato, e a produção de fake news.
O comando solicitava um texto abordando dois tópicos : o tensionamento entre a defesa dos Direitos Humanos e a produção de notícias falsas, tendo como base a vereadora Marielle – a defesa dos Direitos Humanos realizada por ela e as notícias falsas divulgadas após a sua morte.
As fake news também \textbf{estavam} presentes no segundo tópico de análise : os prejuízos que provocam para a sociedade democrática.
Por um lado, os textos selecionados para constituir o enunciado da questão continham apoio suficiente para formulação e articulação de argumentos para discorrer sobre o tensionamento entre a atuação da Marielle em defesa dos Direitos Humanos e fake news.
Por outro, havia ênfase nos aspectos biográficos e na trajetória da vereadora Marielle, o que levou a respostas pouco focadas no comando, que discutiam a vida e a morte da vereadora, sua importância como um ícone dos Direitos Humanos, opiniões em \textbf{torno} do seu assassinato, até mesmo sobre a responsabilidade pelo crime, muitos textos com a presença/defesa de fake news ou produção de novas.
O foco sobre Marielle também ensejou respostas basicamente centradas na questão da violência.
Foram frequentes as interpretações com viés político-partidário, não previstas no padrão de resposta.
O item do comando que solicitava a discussão sobre democracia e fake news encontrou menos apoio nos textos.
Apesar de se esperar que graduandos tivessem maturidade, inclusive política, para demonstrar alguma compreensão dos aspectos principais que envolvem o convívio democrático, esse aspecto foi pouco discutido e, quando presente, quase sempre foi apenas mencionado, numa quase reprodução do comando.
A falta de entendimento do que é o conceito de Direitos Humanos também foi percebida nas respostas, seja para recomendá -lo como importante para a sociedade, seja para criticá -lo como algo que defende bandidos.
Foi marcante a ideia que vincula Direitos Humanos a direitos apenas das minorias.
Não foi incomum encontrar respostas que personificavam os Direitos Humanos, tratando -os como sujeito de alguma ação que deveria ser implementada na sociedade em geral ou em casos específicos.
As fake news foram citadas muitas vezes, mas pouco explicadas, conforme previsto no padrão de resposta.
O argumento utilizado geralmente vinculava as consequências de fake news para a imagem da Marielle, ou seja, centrava no fato em si, já indicado no enunciado, sem qualquer análise ou aprofundamento próprio.
É importante salientar que muitas respostas, provavelmente de grupos vinculados a cursos específicos, continham argumentações recorrentes : os Direitos Humanos como algo vinculado à Constituição e a tratados dos quais o Brasil é signatário ; as fake news como produto do mundo digital e globalizado, relacionado à era da pós-verdade ; a discussão sobre o papel das mídias sociais e da imprensa. produto do mundo digital e globalizado, relacionado à era da pós-verdade ; a discussão sobre o papel das mídias sociais e da imprensa.
Como já apontado, houve baixa quantidade de respostas contemplando o segundo tópico do comando – prejuízos da produção de notícias falsas para a sociedade democrática ; fato motivado, possivelmente, pela falta de um elemento que ressaltasse esse ponto nos textos motivadores.
De modo geral, os estudantes não conseguiram recorrer a conhecimentos tácitos sobre democracia para construir argumentos próprios.
As respostas que abordavam o tema, geralmente \textbf{estavam} vinculadas à ideia do direito à verdade, do direito à integridade da pessoa ou dos direitos gerais garantidos pela Constituição.
A relação da questão com o momento político brasileiro foi indicada em um número expressivo de respostas.
Por vezes, apenas questionando a temática da questão – caso Marielle – em uma prova para avaliar o Ensino Superior no Brasil ou a importância dos Direitos Humanos na atual conjuntura política do país, levando à anulação da resposta.
Houve também interpretações das consequências do uso das fake news no processo eleitoral de 2018, influenciando a escolha consciente e livre das pessoas.
As respostas demonstraram atenção político social frente a esses temas, porém com pouco domínio conceitual do campo de conhecimento sobre Direitos Humanos, fake news e democracia.
Em termos da linguagem adotada pelos alunos, há dois grandes grupos de respostas.
No primeiro grupo, os que responderam de forma clara, com autonomia de pensamento e explicitação de justificativas para os aspectos solicitados no enunciado, demonstrando conhecimento sobre o que significam Direitos Humanos, fake news e democracia, além da articulação entre essas questões.
Por outro lado, um segundo grupo dos que não conseguiram ir além do que foi citado no enunciado, apenas repetindo -o, muitas vezes desviando o foco da questão para temas correlatos, como a vida de Marielle ou a violência.
Nesse segundo grupo, mais facilmente foram encontradas respostas desconexas e mal articuladas, com textos confusos que explicitam o não conhecimento sobre o tema e a falta de adequação à linguagem, necessária a um concluinte do Ensino Superior.
As abordagens dos Direitos Humanos mais recorrentes nas respostas corretas foram : a referência da Marielle como representante/ícone dos DH ; e os DH como direitos previstos na Constituição e em tratados internacionais dos quais o Brasil é signatário.
No entanto, também foram encontradas críticas aos Direitos Humanos, tais como : DH vistos como “ direitos de bandidos ” ; a defesa pouco embasada de frases de efeito como “ direitos humanos para humanos direitos ”.
Também chamou atenção o tratamento dos DH como um personagem, ou um órgão público, evidenciando desconhecimento sobre o tema.
Em relação às fake news, dois argumentos previstos pelo padrão de respostas foram os mais presentes : a noção de que o mundo globalizado e a facilidade de reprodução de informações pelas mídias atuais levam à criação de mentiras que são compartilhadas e que geram consequências ; a Marielle como alvo de fake news que tentaram desconstruir sua luta e seu legado.
Também foi frequente a ideia de que a produção de informações falsas é realizada por uma imprensa tendenciosa com intenção de manipular a informação.
No item relacionado a fake news e democracia, as respostas corretas afirmavam que as fake news inviabilizam o acesso à verdade e à garantia da integridade de um estado Democrático de Direito.
Também foi abordada a influência das fake news nos resultados das eleições dos últimos anos, nos estados Unidos e no Brasil, gerando consequências para a vida das pessoas.
Por fim, um outro aspecto abordado nas respostas a este tópico foi a polarização da sociedade em termos políticos, com a afirmação de que as fake news dificultam o debate, estimulam posições extremas e rivalidades prejudiciais à convivência em uma democracia plena.
A notas mais fracas ( de zero a 35 ) foram cerca de 20 \% das respostas corrigidas – as que não \textbf{estavam} em branco nem foram desconsideradas ou anuladas.
Foram respostas que se restringiam a cópia ou paráfrases do enunciado ; as que tinham foco em desdobramentos políticos e policiais do caso Marielle ; as que tinham argumentação confusa e/ou incompleta.
A grande maioria das notas, em \textbf{torno} de 65 \%, foi mediana, entre 40 a 70.
Nesse grupo ficaram as respostas que continham alguma abordagem adequada sobre Direitos Humanos, sobre as atividades da vereadora Marielle e sobre as fake news divulgadas após sua morte.
As notas mais altas deste grupo foram para respostas que incluíram a questão do tensionamento, quase sempre, em um texto coerente e consistente.
Muitas vezes, a nota mediana foi consequência do baixo índice de respostas à segunda parte do comando, que solicitava a discussão sobre os prejuízos das fake news para a democracia.
As melhores notas só foram obtidas por aqueles que incluíram a discussão sobre os prejuízos das fake news para a democracia.
Apenas 15 \% das respostas corrigidas receberam avaliação igual ou superior a 75.
A deficiência principal detectada pelos corretores foi a dificuldade de grande parte dos alunos para interpretar o enunciado da questão e para se expressar em textos consistentes.
De modo geral, os temas foram tratados de forma fragmentada e pouco analítica.
Além disso, chamou à atenção a predominância do lugar comum na construção argumentativa das respostas, a superficialidade do debate em \textbf{torno} da democracia e o desconhecimento de conceitos como Direitos Humanos, esses confundidos com Instituições e até pessoas. 6.3.1.3 Análise de Conteúdo da Questão Discursiva 2 do Componente de Formação Geral Na Tabela 6.15, mostra -se que o desempenho médio dos estudantes na questão discursiva 2 ( média 30,6 ) foi inferior ao obtido na questão discursiva 1 ( média 31,2 ).
A região Norte foi aquela cuja Média, também nessa questão, foi maior ( 34,6 ), enquanto a de menor Média foi a região Centro-Oeste ( 25,6 ).
Quanto à variabilidade das notas, o Desvio padrão de todo o Brasil foi 23,3, superior ao obtido na questão discursiva 1 ( 23,0 ).
O maior desvio nessa questão foi obtido na região Norte ( 25,6 ), enquanto o menor foi obtido na região Sul ( 21,9 ).
A Mediana de todo o Brasil foi 30,0, a mesma de três regiões ( Nordeste, Sudeste e Sul ).
As exceções foram as regiões Norte e Centro-Oeste, que obtiveram medianas 35,0 e 25,0, respectivamente.
As notas máxima ( 100,0 ) e mínima ( 0,0 ) foram as mesmas em todas as regiões, sem exceção.
No Gráfico 6.8, mostra -se a distribuição das notas na questão discursiva 2 do Componente de Formação Geral.
O intervalo modal foi o ( 30 ; 40 ], com 18,4 \% dos participantes, e seguido de perto pelo intervalo ( 40 ; 50 ], com 16,2 \%.
Destaca -se que quase 14 \% dos alunos deixaram a questão em branco e um pouco mais de 8 \% obteve zero como nota nessa questão.
Nota -se ainda que as notas ficaram mais dispersas em comparação às da questão discursiva de número 1, o que pode ser constatado, também, pela comparação do desvio padrão das notas da questão discursiva 2 ( 23,3 ) e o das notas da questão discursiva 1 ( 23,0 ). 6.3.1.4 Comentários sobre a correção de Conteúdo das respostas à Questão Discursiva 2 Os textos e o comando da Questão 2 de Formação Geral são perfeitamente adequados ao que se espera avaliar em Formação Geral, ao tratar das funções sociais dos museus, contextualizando a temática sob o ponto de vista da memória e da perspectiva de futuro.
Na questão, \textbf{estavam} presentes, principalmente, o conteúdo humanista e crítico, a interpretação e a produção de textos sobre temas como cultura e arte, educação, ciência e democracia, previstos na Portaria que dispõe sobre o componente de Formação Geral do Enade 2018.
Os conhecimentos exigidos para responder ao comando, explicitados no padrão de resposta, vão além do senso comum, por se exigir a compreensão de conceitos como “ função social ” ( de Instituições sociais ), “ memória ” ( no sentido de memória social ), “ perspectiva de futuro ” ( no contexto sócio-histórico-cultural ) e “ sociedade contemporânea ” ( que deveria independer da interpretação de \textbf{cada} grupo social a respeito da sociedade na qual vive ).
No entanto, apesar de envolver tais conceitos, espera -se que alunos concluintes de cursos superiores de graduação sejam capazes de discutir funções de Instituições sociais, especialmente porque o enunciado trazia fato recente e de grande repercussão como apoio.
O padrão de resposta elencou as funções sociais que poderiam ser abordadas nas respostas dos estudantes, como também argumentos que poderiam ser usados ao longo das mesmas.
A tarefa dos estudantes era a construção de um texto expositivo que articulasse a importância de museus, suas funções sociais e as perspectivas de memória e de futuro, o que se mostrou de dificuldade de média a difícil.
Apesar de o enunciado ser claro, muitos estudantes, provavelmente influenciados pelos textos, limitaram -se a discutir aspectos relacionados ao Museu Nacional e, em particular, ao incêndio que o destruiu.
Foram poucos os que interpretaram corretamente o escopo do comando, que solicitava a discussão das funções sociais de museus de uma forma geral, no sentido de argumentar a importância deste tipo de Instituição.
As quatro funções sociais elencadas no padrão de resposta foram encontradas na quase totalidade das redações dos estudantes de forma articulada com a importância dos museus.
Entretanto, quase nunca os textos abordavam mais do que uma função social e vinculavam, de forma clara, a importância desse tipo de Instituição para a sociedade contemporânea.
Além disso, a solicitação de contextualização sob o foco museu/memória/perspectiva de futuro parece não ter sido bem compreendida e, em consequência, \textbf{esteve} pouco presente nas respostas.
Também se observou confusão entre os conceitos de história e memória e uma visão estreita do significado de sociedade contemporânea, às vezes, restrita à cidade do respondente, outras, ao Rio de Janeiro, ou ainda a um tipo de sociedade idealizada.
A relevância que os textos de apoio deram ao caso particular do incêndio do Museu Nacional do Rio de Janeiro e aos prejuízos causados ao seu acervo, induziu a maioria dos estudantes a elaborar respostas tratando das funções a partir da perda e, em particular, à perda do tipo de acervo daquele museu e do tipo de pesquisa que lá se desenvolve, divergindo do comando da questão, que solicitava texto abordando funções sociais de museus em geral.
Também foram encontradas muitas respostas que se limitavam a discutir o que o estudante considerava como possíveis causas do incêndio, a falta de manutenção, questões políticas e de gestão, sem responder ao que foi solicitado.
Foi também recorrente a transcrição de frases dos textos do enunciado nas respostas, diretamente ou por paráfrases, revelando \textbf{certa} dificuldade de criação textual.
A análise das respostas dos estudantes, além de revelar dificuldade de abordar a temática esperada, revelou outros dois aspectos predominantes.
Um deles, como já comentado, foi a ênfase no incêndio do Museu Nacional, que orientou grande parte dos textos dos estudantes.
Outro, também evidenciando desvio do esperado, foi o foco na perda patrimonial.
Nesse grupo, houve respostas que até apresentaram uma função social, mas olhada pelo ponto de vista da perda.
Nessa linha, muitas respostas enalteciam a importância do Museu Nacional, sua relevância no cenário nacional, o impacto da perda de seu acervo, isso complementado por frágeis reflexões sobre as funções sociais.
Em consequência, a perspectiva de desdobramentos mais ricos, que tratassem de aspectos sociais relacionados a diferentes tipos de museu foi prejudicada.
Muitas respostas trataram da preservação da memória e/ou guarda do acervo ; identidade, tanto individual quanto coletiva ; conhecimento sobre antepassados ; evolução da sociedade e dos seus costumes ; valorização e disseminação do conhecimento por meio da liberação do acesso aos acervos ; bem como atividades pedagógicas e exposições de arte.
O aspecto produção de pesquisa e conhecimento, embora frequente, foi mencionado de forma muito simples, sem qualquer aprofundamento, para além do que já era mencionado no enunciado.
Frequentemente apresentava -se o museu como produtor de pesquisas para o campo da saúde, o que evidencia desconhecimento das reais possibilidades de pesquisa, mesmo no Museu Nacional.
Quanto a perspectiva de futuro, quando presente no texto, era mencionada de forma breve, com afirmações genéricas sobre a possibilidade de melhoria do futuro pelo simples conhecimento do passado, numa percepção de que o passado determina o futuro, de forma linear.
Foram raras as respostas com alguma argumentação que associasse pesquisa, memória, identidade e reflexão social como base para a construção de alternativas de futuro.
A relação entre turismo e museu apareceu considerando a visita a museus como forma de lazer e, algumas vezes, como geradora de recursos.
O desempenho pode ser considerado fraco.
Muitas respostas deixavam claro que o estudante nunca tinha visitado um museu e, de uma forma geral, observou -se desconhecimento sobre suas funções sociais.
Não se esperava que história, memória social e museologia fossem assuntos dominados pelos estudantes das diversas áreas que participaram do Enade 2018, mas esperava -se que essa temática pudesse ser abordada do ponto de vista das experiências pessoais ou de conhecimento social.
Como uma vivência própria com os museus também se revelou frágil ou inexistente, o conhecimento sobre o assunto ficou bem restrito aos textos motivadores e ao senso comum, construído, basicamente, pelo noticiário em \textbf{torno} do incêndio.
O universo dos argumentos dos estudantes se mostrou muito limitado. \textbf{Cada} dimensão foi praticamente representada por uma ideia simplista.
De forma recorrente : a preservação da memória para construção da identidade de um povo ; a compreensão do passado por meio das pesquisas como oportunidade para melhorar o futuro ; o museu como um espaço em que há um aprendizado estimulado pelo lazer.
Houve, também, uma grande concentração de respostas associando o museu a um local de guarda de acervo e de história.
As respostas à Questão 2 evidenciaram dificuldades conceituais.
Além dos conceitos de história, memória e museu se misturarem, parecendo não se distinguirem, ao mencionarem a pesquisa, também se observou falta de distinção entre pesquisas colegiais, com objetivos de aprendizado, e as realizadas por pesquisadores, com objetivos de avançar cientificamente.
A falta de compreensão do que se pedia no comando da questão e a dificuldade de expressão linguística foram dois aspectos marcantes na avaliação dos corretores.
Foi frequente os textos apresentarem ideias soltas ou de forma confusa, sem constituir unidade.
O vocabulário usado se revelou bem restrito, e fortemente apoiado no coloquial, na linguagem oral.
Predominou o senso comum em lugar de uma reflexão mais formal envolvendo os conceitos envolvidos. 6.3.1.5 Análise de Língua Portuguesa das Questões Discursivas do Componente de Formação Geral Os dados de Serviço Social, obtidos a partir das respostas às questões discursivas do Componente de Formação Geral, no que tange à Língua Portuguesa, encontram -se na Tabela 6.16 e no Gráfico 6.9.
Nesse aspecto, os alunos de todo o Brasil obtiveram Média 51,2.
A maior Média, em relação à Língua Portuguesa, foi obtida na região Sudeste ( 51,8 ), e a menor, na região Sul ( 51,0 ).
Quanto à variabilidade das notas, o Desvio padrão de todo o Brasil foi 23,0.
O menor Desvio padrão foi obtido na região Sul ( 21,0 ), e o maior Desvio padrão foi obtido na região Norte ( 24,1 ).
A Mediana das notas de Língua Portuguesa foi 57,5 para o Brasil, como um todo, a mesma obtida nas regiões Norte e Sul.
A Mediana das regiões Nordeste e Sudeste foi 60,0.
A menor Mediana foi na região Centro-Oeste ( 55,0 ).
A nota Máxima para todo o Brasil foi 92,5, com, pelo menos, um aluno tirando essa nota na região Sudeste.
Nas regiões Norte, Nordeste e Sul, a nota Máxima foi 90,0, e na Centro-Oeste, foi 87,5.
Já a nota Mínima foi zero em todas as regiões do país.
No Gráfico 6.9, mostra -se a distribuição das notas de Língua Portuguesa do Componente de Formação Geral.
Observa -se que a maior frequência ( cerca de 27 \% ) corresponde à dos alunos que obtiveram nota no intervalo ( 60 ; 70 ].
Destacam -se, também, os alunos que deixaram ambas as questões em branco, representam 9 \% do total, caracterizando -se como máximo local. 6.3.1.6 Comentários sobre a correção das respostas de Formação Geral com respeito à Língua Portuguesa As duas questões discursivas de Formação Geral ensejaram duas configurações textuais distintas, tendo em vista os seus encaminhamentos : a Questão 1 teve um caráter mais opinativo e a Questão 2, mais expositivo, conforme se explicita a seguir.
Ambas as questões permitiram a observação do desempenho linguístico dos participantes em situação formal, competência exigida no processo de formação universitária.
As habilidades de produção textual envolvidas têm, igualmente, como condição básica, o domínio do padrão formal culto da Língua Portuguesa, no que diz respeito aos aspectos relativos às questões de caráter ortográfico, textual, morfossintático e vocabular.
O encaminhamento das questões evidencia a expectativa de que o participante utilizasse seus conhecimentos sobre os dois temas e estruturasse seus textos de acordo com as características do registro formal adequado à situação comunicativa.
Tal configuração determina exigências quanto : à adequação da seleção vocabular, ao desenvolvimento do conteúdo, à estruturação sintática dos períodos, à organização lógica das ideias, à utilização de procedimentos de encadeamento textual e de referenciação, à obediência às exigências morfossintáticas próprias da modalidade escrita da norma-padrão, ao respeito às regras ortográficas e às regras de acentuação gráfica.
Esse encaminhamento exigiu do graduando o domínio do modo de organização textual de base dissertativo-argumentativa, essencial ao processo de formação universitária.
As habilidades de produção textual envolvidas têm como condição básica o domínio do padrão formal culto da Língua Portuguesa, no que diz respeito aos aspectos relativos às questões de caráter ortográfico, textual, morfossintático e vocabular.
O enunciado da Questão 1 tomava como base três fragmentos de textos jornalísticos publicados em sites da internet, que abordavam o aumento do assassinato de minorias e de defensores de direitos humanos no Brasil, entre eles a vereadora Marielle Franco.
O comando solicitava que o participante desenvolvesse um texto dissertativo e argumentativo com as temáticas Fake News, Democracia e Direitos Humanos.
Embora não tenha havido referência explícita à tipologia textual esperada para o desenvolvimento da resposta, essa proposta encaminhou o participante para a elaboração de um texto opinativo, com o objetivo de discutir o problema e manifestar um posicionamento crítico.
O enunciado da Questão 2 abordava o incêndio do Museu Nacional do Rio de Janeiro, ocorrido em 2018, e seus impactos na ciência brasileira e internacional, tomando como base dois fragmentos de textos jornalísticos publicados em sites da internet.
O comando solicitava que o participante desenvolvesse um texto abordando três aspectos da função social dos museus na sociedade contemporânea, encaminhando -o para a elaboração de um texto expositivo que confirmasse a posição defendida nos fragmentos motivadores a respeito da importância dessas instituições.
Para atender a essas exigências, o padrão de resposta utilizado na avaliação das questões discursivas de Formação Geral considerou aspectos relevantes ao bom desempenho linguístico como competências distintas, de modo a permitir um mapeamento detalhado do domínio dos recursos disponíveis na Língua Portuguesa para a comunicação escrita formal : ortográficos, textuais, morfossintáticos e vocabulares ( ver detalhes no Anexo VIII – Padrão de Respostas ).
Aspectos Ortográficos Para avaliar a competência relativa ao domínio das convenções ortográficas da norma-padrão da Língua Portuguesa, observou -se o respeito às regras de acentuação gráfica e da grafia padrão das palavras ( com ausência de abreviaturas próprias da linguagem da internet ), de acordo com as convenções estabelecidas pela legislação em vigor e consubstanciadas no Vocabulário Ortográfico da Língua Portuguesa, editado pela Academia Brasileira de Letras.
Aspectos Textuais Para avaliar a competência relativa ao domínio dos procedimentos de estruturação textual, procurou -se observar a utilização de mecanismos linguístico-discursivos que estabeleçam o encadeamento lógico entre as partes do texto, de forma a garantir a progressão e a coerência textuais.
Esta competência destaca -se por sua relevância, já que são avaliadas a clareza e a unidade textual.
Aspectos Morfossintáticos Para avaliar o domínio dos diferentes aspectos morfossintáticos próprios da modalidade escrita formal da norma-padrão da Língua Portuguesa, observou -se se o texto produzido atendeu às seguintes exigências : a concordância nominal, a concordância verbal, a regência nominal, a regência verbal, a flexão nominal, a flexão verbal, a correlação entre os tempos verbais, a colocação pronominal e a utilização de sinais de pontuação que contribuam para a organização lógica da frase e do texto.
Aspectos Vocabulares Para avaliar a adequação da seleção vocabular à modalidade escrita formal da Língua Portuguesa, exigida pela situação comunicativa, observou -se o atendimento às seguintes exigências : precisão na seleção/utilização do vocabulário relacionado à temática solicitada pela questão ; ausência de marcas da oralidade, como termos de sentido muito genérico ( “ coisa ”, “ negócio ”, “ você ” ) e termos de registros mais informais ( como gírias, jargões, frases feitas, ditados populares, termos regionais ).
Assim, espera -se que o participante respeite a adequação vocabular não usando gírias ou expressões coloquiais, evite repetição desnecessária de palavras e utilize um vocabulário mais formal, como solicitado por um texto dissertativo.
Para efeito de pontuação, a avaliação do desempenho linguístico considerou três grandes grupos de competências, agrupando aspectos explicitados, com os seguintes pesos relativos : aspectos ortográficos ( 20 \% ) ; aspectos textuais ( 40 \% ), aspectos morfossintáticos e vocabulares ( 40 \% ).
Os critérios de pontuação utilizados no padrão de resposta procuraram atender à diversidade de desempenho dos graduandos.
Nos textos analisados, observaram -se, assim, vários níveis de atendimento às exigências do padrão formal próprio da modalidade escrita da Língua Portuguesa, configurando um"continuum".
Por um lado, textos que traduzem com exatidão a expectativa de uso formal da língua, em todos os níveis analisados, marcados por clareza e unidade ; e, no outro extremo, textos quase incompreensíveis devido ao grande número de desvios de todos os tipos, principalmente os de caráter textual.
As duas questões discursivas de Formação Geral analisadas revelaram desempenhos distintos dos participantes.
A Questão 1 propiciou a oportunidade de maior desenvolvimento do tema solicitado, gerando consequências na estruturação textual, já que os textos foram mais longos, mais elaborados e fluentes.
Em virtude da polêmica sobre a existência das fake news e sobre a visão pública da personagem da vereadora Marielle Franco, observaram -se respostas que polemizaram o tema, favorecendo a construção argumentativa do texto.
A Questão 2 teve um comportamento distinto, devido ao caráter mais expositivo do enunciado e ao reduzido conhecimento dos participantes sobre o tema.
Assim, os textos foram mais curtos, frequentemente com repetição de ideias dos textos motivadores.
Quanto aos aspectos linguísticos analisados durante esta avaliação, observaram -se os seguintes resultados : Aspectos ortográficos : O desempenho dos participantes revelou uma diferença muito grande nos dois aspectos analisados nesta competência : baixo índice de desvios da grafia padrão e grande índice de desvios de acentuação.
Em vários casos, ocorre ausência completa de acentuação gráfica em todas as palavras do texto.
Os resultados revelaram que a tendência dominante entre os universitários brasileiros é a eliminação da acentuação gráfica, talvez motivada pelos hábitos relacionados às redes sociais e pela ausência de esclarecimento dos meios de comunicação, das autoridades e das escolas sobre as decisões do Acordo Ortográfico da Língua Portuguesa de 1990.
Os casos mais sistemáticos de eliminação do acento indicador da sílaba tônica foram : - palavras proparoxítonas ( por exemplo “ generos ”, “ arbitro ”, “ politico ”, “ publicas ”, “ numero ” ) ; - palavras paroxítonas terminadas em ditongo crescente ( por exemplo “ area ”, “ noticiarios ”, “ individuo ”, “ historia ”, “ varios ”, “ noticia ”, “ negligencia ”, “ ocorrencia ” ) ; - palavras oxítonas ( por exemplo “ ninguem ”, “ esta ”, “ ate ”, “ tambem ” ).
Por outro lado, destaca -se o uso indevido do acento gráfico em palavras como “ melâncolia ” e “ intervênção ”, por exemplo.
Quanto ao domínio das convenções relativas à grafia das palavras, observaram -se desvios esporádicos como : hipercorreção pela escolha de “ e ” no lugar de “ i ”, por influência de hábitos da oralidade ( “ descriminação ” no lugar de “ discriminação ”, “ entervenção ” no lugar de “ intervenção ” ) ; desvios de grafia relacionados à variação diastrática, como por exemplo “ perca ”, “ subjulgadas ”, “ precoseito ”, “ soubre ”, “ vecendo ”, “ indesencia ”, “ apolojia ”, “ fulga ”, “ dereitos ” ;...
Observaram -se, também, muitos casos de inadequação no uso da maiúscula : uso indevido para destacar determinadas palavras-chave do texto, como “ Brasileiros ”, “ Homens ” ; ausência de maiúscula para grafar nomes próprios, como"marielle".
Expressões como"Direitos Humanos"foram grafadas diversificadamente, com maiúsculas e com minúsculas.
Destaque -se, também, a maioria dos participantes que grafam os textos inteiramente em caixa alta, o que impede a identificação de desvios relacionados ao uso de maiúsculas.
Vale observar, também, que, ao contrário do que se esperava, não apareceram abreviaturas próprias do “ internetês ”, relacionadas ao uso de redes sociais e e-mails, como por exemplo"vc".
Aspectos textuais : Para alguns dos graduandos essa competência se revelou como a mais problemática, tendo em vista os inúmeros problemas observados, desvios acumulados durante toda a formação escolar.
São eles : sequência justaposta de ideias sem encaixamentos sintáticos, provocando uma redução de estruturas subordinadas, ao lado do aumento na frequência de estruturas coordenadas e absolutas ; frases fragmentadas que comprometem a estrutura lógico-gramatical ; frases formadas apenas por oração subordinada, sem oração principal ; redução no uso de conectores para expressar relações lógicas essenciais à construção do texto como consequência da mudança de estruturação frasal ; emprego inadequado do pronome relativo ( com omissão da preposição ou a utilização do pronome relativo inadequado, como “ onde ” ), refletindo hábitos da oralidade ; ausência de recursos de referenciação, como substituição de termos por sinônimos, hiperônimos, nominalizações, expressões metafóricas. lógico-gramatical ; frases formadas apenas por oração subordinada, sem oração principal ; redução no uso de conectores para expressar relações lógicas essenciais à construção do texto como consequência da mudança de estruturação frasal ; emprego inadequado do pronome relativo ( com omissão da preposição ou a utilização do pronome relativo inadequado, como “ onde ” ), refletindo hábitos da oralidade ; ausência de recursos de referenciação, como substituição de termos por sinônimos, hiperônimos, nominalizações, expressões metafóricas.
Esses problemas revelam dificuldades em relação à estrutura formal do texto produzido, o que é preocupante ao se levar em conta que são graduandos em fase final de formação.
Em alguns textos, falta um mínimo de textualidade e de domínio do registro padrão da língua.
Na verdade, observaram -se relações linguísticas quase agramaticais, como as estabelecidas pela sequência de gerúndios sem o apoio de um ponto de partida para a organização das informações gramaticais e semânticas.
Quanto à utilização dos sinais de pontuação, observou -se uma ausência dos mesmos na maioria dos textos analisados.
Foi muito frequente a ocorrência de parágrafos sem marca interna de pontuação para separar os períodos.
São os seguintes os tipos de problemas mais encontrados :  vírgula : utilização de vírgula para separar o sujeito e o predicado ; uso de vírgula no lugar do ponto para separar ideias que constituem períodos distintos ; ausência de vírgula para separar oração adjetiva explicativa ;  ponto e vírgula : utilização do ponto e vírgula no lugar de vírgula ;  ponto final : ausência de ponto final para separar períodos.
Aspectos morfossintáticos e vocabulares : Os resultados são transparentes em relação aos aspectos mais problemáticos no desempenho dos participantes nestas duas competências.
Em relação aos aspectos morfossintáticos, seguem algumas observações sobre os desvios mais frequentes.
Quanto à regência, o desvio mais frequente foi a falta do sinal indicativo da crase – isso revela que o usuário não tem consciência de que, sob a forma do termo “ a ”, existe a presença de uma contração entre a preposição “ a ” ( exigida pela regência do termo anterior ) e o artigo definido “ a ”.
Outro problema relacionado à regência verbal e à nominal, encontrado frequentemente nas respostas, foi a ausência de preposição antes de pronome relativo, processo generalizado na modalidade oral da língua, em situações de registro informal.
Apesar da possibilidade de que essa alteração de regência se generalize no padrão escrito da Língua Portuguesa, como já \textbf{está} ocorrendo até em textos jornalísticos, o não emprego da preposição foi considerado como desvio neste processo de avaliação.
Outro desvio de regência significativo foi a utilização inadequada de uma preposição ou sua ausência após o verbo ou o nome ( substantivo ou adjetivo ). na modalidade oral da língua, em situações de registro informal.
Apesar da possibilidade de que essa alteração de regência se generalize no padrão escrito da Língua Portuguesa, como já \textbf{está} ocorrendo até em textos jornalísticos, o não emprego da preposição foi considerado como desvio neste processo de avaliação.
Outro desvio de regência significativo foi a utilização inadequada de uma preposição ou sua ausência após o verbo ou o nome ( substantivo ou adjetivo ).
A concordância verbal e a concordância nominal apresentaram alguns desvios muito frequentes.
Quanto à concordância de número, observou -se, como apontam as pesquisas nessa área, ausência de marca ( com sujeito anteposto ou posposto ) ou uso indevido ( uso inadequado da marca de plural comandado pelo núcleo plural da locução adjetiva, apesar de o substantivo que funciona como núcleo do sintagma nominal \textbf{estar} no singular ).
Uma ocorrência que se destacou foi a ausência de acento circunflexo na forma plural do presente do indicativo do verbo “ ter ”, que foi considerada como um desvio na concordância verbal e não na acentuação gráfica.
Quanto à concordância de gênero, vários casos foram observados, normalmente no âmbito de sintagmas nominais longos, em que o adjetivo \textbf{está} afastado do substantivo.
Quanto à questão da colocação pronominal, foram poucos os desvios observados.
Concluiu -se que, no registro escrito formal, a maioria dos participantes já incorporou regras como a não introdução da frase por um pronome oblíquo e a próclise na presença de um termo atrator.
Várias marcas de oralidade foram identificadas, embora não com alta frequência : uso do pronome relativo “ onde ” como relativo universal, falta de artigo definido antes de substantivo isolado ou antes de termos em enumerações.
Quanto aos aspectos vocabulares, alguns tipos de inadequação foram observados : expressões da oralidade ; seleção vocabular incompatível com o contexto, gerando falta de inteligibilidade ; falta de domínio de vocabulário mais abstrato e de maior complexidade, essencial ao desenvolvimento do texto de base dissertativa.
O principal aspecto observado foi a excessiva repetição de \textbf{certas} palavras, como o termo"pessoa", por exemplo, revelando limitação de repertório vocabular.
Em função do tema solicitado na Questão 1, diferentes grafias da expressão fake news não foram consideradas.
Para sintetizar, os problemas que mais se destacaram nos textos analisados foram : a ) truncamentos sintáticos, caracterizados por estruturas frasais incompreensíveis devido à ruptura da complexidade sintática própria da modalidade escrita ; b ) redução de estruturas subordinadas, compensada pelo aumento na frequência de estruturas coordenadas e absolutas, por um lado, ou pela elaboração de estruturas textuais prejudicadas pelo excesso de ideias sem a devida conexão subordinativa ; c ) redução no uso de conectores para expressar relações lógicas essenciais à construção do texto, substituídos pela exigência de inferência por parte do interlocutor para suprir a sua ausência ; d ) redução no uso do subjuntivo, ao lado da ampliação do uso do indicativo, combinado a estruturas frasais coordenadas ou absolutas ; d ) empobrecimento do processo de referenciação, com a repetição de palavras, sem processos mais sofisticados de substituição ; e ) simplificação da marcação da categoria tempo na morfologia verbal ; f ) desrespeito às regras de concordância verbal e nominal próprias do padrão formal ; g ) ausência do acento grave indicador da crase, revelando falta de conhecimento dos critérios de regência verbal e nominal ; h ) falta de domínio de vocabulário mais abstrato e de maior complexidade, essencial ao desenvolvimento do processo dissertativo ; i ) redução do emprego da acentuação gráfica.
Como resultado da correção, merece atenção a constatação de que a modalidade escrita tem apresentado a tendência a uma evidente simplificação, aproximando -se das características da modalidade oral da Língua Portuguesa.
No caso dos textos de base dissertativa ( expositiva ) e de base argumentativa, inscritos em um registro formal, a distância entre as duas modalidades é ainda maior, o que provoca desvios recorrentes em todos os quatro aspectos analisados : ortográficos, textuais, morfossintáticos e vocabulares.
Destaca -se, nessa avaliação, o comprometimento textual, por meio de estruturas fragmentadas e/ou truncadas, rompendo a complexidade sintática esperada no padrão formal. 6.3.2 Componente de Conhecimento Específico Na parte da prova relativa às questões discursivas do Componente de Conhecimento Específico ( Tabela 6.17 ), observa -se que a Média foi bem mais baixa do que para as questões discursivas do Componente de Formação Geral.
Enquanto no Componente de Formação Geral a Média para estudantes de Serviço Social de todo o Brasil foi 34,8, na parte de Conhecimento Específico, a Média foi 21,7.
A maior Média deste componente foi obtida pelos estudantes da região Nordeste ( 25,7 ), e a menor, pelos da região Sul ( 18,9 ).
Quanto à Quanto à variabilidade das notas, o Desvio padrão de todo o Brasil foi 18,4.
O maior Desvio padrão foi encontrado na região Nordeste ( 20,1 ), e o menor, na região Sul ( 15,7 ).
A maior nota Máxima, 90,0, foi obtida, também, nas regiões Norte, Sudeste e Sul.
Nas demais regiões, as notas Máximas foram 83,3, na região Nordeste e 86,7, na Centro-Oeste.
A nota Mínima ( 0,0 ) foi obtida por, pelo menos, um aluno em todas as regiões do Brasil.
A Mediana do Brasil, como um todo, foi 20,0.
A maior Mediana foi obtida nas regiões Norte e Nordeste ( 23,3 ), e a menor, nas regiões Sudeste, Sul e Centro-Oeste ( 16,7 ).
No Gráfico 6.10, representa -se a distribuição das notas nas questões discursivas no Componente de Conhecimento Específico.
A moda dessa distribuição ocorre no intervalo [ 0 ; 10 ], com mais do que 36 \% do total de participantes.
Observa -se que nesses intervalos \textbf{estão} computadas as ocorrências de respostas em branco para as três questões discursivas de Conhecimentos Específicos.
A análise de \textbf{cada} uma destas questões será feita a seguir. 6.3.2.1 Análise da Questão Discursiva 3 do Componente de Conhecimento Específico Na questão 3, cujos resultados aferidos encontram -se descritos na Tabela 6.18, a Média dos estudantes de todo o Brasil foi 25,0, e a questão de desempenho mediano dentre as três discursivas de Conhecimentos Específicos.
A menor Média nessa questão foi obtida pelos alunos da região Centro-Oeste ( 21,7 ), enquanto a maior Média foi obtida na região Nordeste ( 29,5 ).
Quanto à variabilidade das notas, o Desvio padrão de todo o Brasil foi 24,2.
O maior Desvio padrão foi obtido na região Norte ( 25,7 ), enquanto o menor foi obtido na região Sul ( 22,5 ).
A nota Máxima, 100,0 pontos, foi alcançada por, pelo menos, um aluno de todas as regiões, com exceção da região Centro-Oeste, que teve nota Máxima 90,0.
A Mediana do Brasil, como um todo, foi 30,0, a mesma para a maioria das regiões, exceto nas regiões Sudeste e Centro-Oeste que foi 20,0.
A nota Mínima ( 0,0 ) também foi a mesma em todas as regiões do Brasil.
No Gráfico 6.11, mostra -se a distribuição das notas na questão discursiva 3 do Componente de Conhecimento Específico da área de Serviço Social.
Sem se levar em consideração que próximo de 15 \% dos estudantes deixaram a questão em branco e que cerca de 23 \% receberam nota zero, a moda dessa distribuição seria o intervalo ( 20 ; 30 ], seguido do intervalo ( 40 ; 50 ]. 6.3.2.2 - Comentários sobre as respostas à Questão Discursiva 3 A Questão 3 abordou dois eixos principais : a relação entre movimentos sociais e democracia ; e o caráter seletivo da ação penal na \textbf{criminalização} das lutas e movimentos sociais.
Para subsidiar a questão, foram apresentados dois trechos de textos abordando o tema, os quais sinalizavam a expectativa de resposta a ser dada pelos estudantes, tendo em vista a sua função de introduzir o tratamento dos eixos mencionados acima.
A Questão 3 requeria o desenvolvimento de um conteúdo que articula as discussões teóricas referentes à relação entre o estado e a sociedade civil, mediada pelos movimentos sociais, fundamental no âmbito dos Cursos de Graduação em Serviço Social.
A questão, também, exigia a compreensão de aspectos históricos sobre a formação social brasileira, priorizando a particularidade da constituição da democracia frente aos traços de uma cultura política autoritária e repressiva.
Apontava, ainda, a necessidade de abordagem dos movimentos e lutas sociais protagonizados por diferentes sujeitos coletivos que atuam na consolidação dos direitos sociais e humanos.
Delimitava, desse modo, conteúdos presentes nas Diretrizes Curriculares da área de Serviço Social, assim como contemplava o esforço de reflexão e a competência necessária ao estabelecimento de mediações teóricas e históricas para a formulação da resposta à questão.
Era, portanto, uma questão que solicitava dos estudantes a compreensão de conteúdos de grande relevância na formação profissional em Serviço Social, tais como : a herança colonial, escravista e patriarcal na formação social, política, cultural e econômica do Brasil ; e os movimentos sociais e as lutas contra opressões de classe, de gênero, étnico-raciais, de sexualidade, geracionais, dentre outras.
O profissional de Serviço Social atua no enfrentamento da questão social mobilizando políticas sociais públicas e privadas, bem como respostas oriundas das organizações da sociedade civil e dos movimentos sociais.
Para tanto, exige -se desse profissional uma formação crítica, comprometida com valores éticos e políticos estabelecidos nas normativas das organizações profissionais, como Código de Ética Profissional e Lei de Regulamentação da Profissão.
Em sua formação acadêmica são contemplados conteúdos que fomentam o desenvolvimento de habilidades e competências teóricas, metodológicas, técnicas, operativas, éticas e políticas em três direções : a análise do movimento histórico e das particularidades do desenvolvimento do capitalismo na sociedade brasileira ; a compreensão do significado social da profissão e de suas contradições ; e a identificação de demandas da realidade social articuladas à questão social com vistas à formulação de respostas profissionais.
Em consonância com essas diretrizes, a Questão 3 atendia plenamente aos requisitos acadêmicos e interventivos esperados do profissional de Serviço Social.
Como sinalizado, a questão formulada exigiu a explicitação de conhecimentos sobre a realidade sócio-histórica brasileira no que tange ao processo de amadurecimento de sua democracia, identificando -se rupturas e potencialidades advindas das organizações de classe e populares ao longo da história.
Propunha, ainda, o enquadramento desse processo nos limites da sociedade capitalista, que coloca em tensão projetos de sociedade, divergentes e antagônicos, e gera padrões de controle e regulação na relação estado e sociedade, de caráter seletivo marcados, não raras vezes, por violência e repressão.
A questão apontava, assim, para a análise dos processos contemporâneos de luta pela democracia ( marcadamente expressos nas pautas antirracistas, anticlassistas, antipatriarcais e antilgbtfóbicas ), sem perder de vista uma compreensão histórica da experiência brasileira em \textbf{torno} das lutas sociais.
O enunciado da questão era claro e apresentava elementos importantes para a organização da resposta, pois os textos do enunciado são de autores que integram o universo cultural dos estudantes de Serviço Social.
Para além de demarcarem as vertentes teóricas, as citações se complementavam no que diz respeito à abordagem da relação entre democracia, direitos e ação dos movimentos sociais, como, também, no que tange às particularidades da cultura política brasileira, os traços autoritários e repressivos do estado e suas bases de justificação.
Destaca -se, inclusive, a exigência de um conteúdo mais \textbf{apurado} do conceito de “ democracia restrita ” em Florestan Fernandes.
A questão pode ser considerada de nível de dificuldade médio, porque exigia uma compreensão das movimentações populares no cenário atual, à luz de aspectos como seletividade, controle/regulação, relação estado e sociedade civil.
Cabe destacar, inclusive, que, em relação às demais questões discursivas da prova, a Questão 3 foi a que trazia, no enunciado, mais elementos teóricos e históricos e a que menos se aproximava de fatos cotidianos amplamente encontrados na grande mídia.
Mas foi, também, a questão que mais introduziu conteúdos centrais dos componentes curriculares da área, e, no que concerne ao grau de profundidade, havia necessidade de conhecimentos de vários conteúdos pertinentes à história, aos movimentos sociais e à democracia no Brasil.
Cabe ressaltar, entretanto, que a menção à ideia de controle penal do estado sobre os movimentos sociais suscitou reflexões direcionadas ao sistema judiciário e, sobretudo, carcerário, tendo em vista a gravidade e evidência desse debate na atualidade.
Como será visto a seguir, algumas respostas evidenciaram a dificuldade de discernimento dos estudantes entre os processos de \textbf{criminalização} das lutas e movimentos sociais e as questões relativas à criminalidade e à segurança pública no país.
O padrão de resposta foi construído de forma adequada ao enunciado da questão. desigualdades como o Brasil.
Não obstante a ausência de referência explícita aos temas do estado, da Sociedade Civil e do modo de produção, muitas respostas trouxeram apontamentos sobre a estrutura da vida social, sustentada na relação entre essas instâncias que compõem a superestrutura e a base econômica e que fundam a compreensão dos fenômenos sociais em \textbf{cada} período histórico.
Considerando -se o grau de profundidade da temática da questão e sua amplitude é possível considerar que o tempo de prova e o espaço reservado para a resposta da questão não tenham sido plenamente suficientes.
Entretanto, ao que tudo indica, a fragilidade no desenvolvimento da resposta \textbf{esteve} mais relacionada a fatores que dizem respeito à formação do profissional, tais como capacidade de leitura e distanciamento no trato dos enunciados, apropriação e domínio categorial, conhecimento histórico da realidade brasileira e competência no estabelecimento de mediações teóricas.
Considerando -se as observações anteriores, no que diz respeito aos conteúdos presentes nas citações e ao que se exigiu no comando, ficou bem evidente, a possibilidade de abordagem da questão com enfoques diversos dos apontados no padrão de resposta.
Desse modo, muitas respostas foram formuladas a partir da relação entre estado, a Sociedade Civil e as particularidades do modo de produção capitalista, interpretação coerente com o padrão de resposta, ainda que nele não se tenha referido às mesmas categorias teóricas e conceitos.
A diversidade de respostas não pode ser justificada em razão de uma suposta incompreensão do enunciado, já que sua formulação foi bem construída, e sua apresentação, em uma única frase e indicando dois itens distintos para a resposta, favoreceu em muito a delimitação objetiva do que \textbf{estava} sendo requerido.
A diversidade pode ser observada nos aspectos abaixo relacionados. - Respostas que parafrasearam as citações.
Muitas respostas não produziram nenhum texto autoral, recorrendo apenas à apropriação das ideias e palavras constantes nas citações, ainda que alterada a \textbf{ordem} entre elas.
Tratava -se, também, em muitas respostas, de uma mera transformação do enunciado da questão no texto de resposta. - As referências ao Projeto Ético e Político da profissão, em sua grande maioria, não auxiliaram na formulação de uma resposta sintonizada com o padrão de resposta.
Ainda que não exclusivamente – mas majoritariamente –, quando esse tema aparecia nas respostas dos estudantes, havia, também, uma distorção na abordagem do papel histórico dos movimentos sociais para a consolidação da democracia e dos direitos sociais para a profissão.
Em muitos casos, adentrou -se no movimento de reconceituação do Serviço Social para abordar as conquistas democráticas no país, ampliando, erroneamente, a contribuição da profissão.
Não se pode, aqui, justificar tal opção como um equívoco de interpretação ou desvio na construção do raciocínio.
O fato é que a incorporação desse tema, constitutivo do Código de Ética Profissional, em relação aos temas da democracia e dos direitos sociais e humanos, pouco contribuiu para a elaboração de respostas corretas.
Além disso, muitas respostas tratavam da relação entre o exercício profissional e os movimentos sociais, destacando a atribuição profissional de prestar assessoria e consultoria a tais movimentos, em total desarticulação com o que foi solicitado. - Majoritariamente, as respostas que apresentavam conteúdos corretos \textbf{estavam} incompletas, ou seja, a primeira parte das respostas se coadunava com o padrão de resposta, já a segunda, não.
Observou -se, deste modo, que os estudantes conseguiam estabelecer os vínculos históricos e teóricos da “ relação entre movimentos sociais e democracia no Brasil ”, mas não obtiveram o mesmo êxito no tocante ao “ caráter seletivo da ação penal do estado na \textbf{criminalização} das lutas e movimentos sociais da sociedade brasileira na atualidade ”.
As respostas incompletas indicavam mais facilidade na apreensão histórica e estrutural da luta de classes no sistema capitalista e do papel dos movimentos sociais do que em tratar as expressões contemporâneas das lutas sociais, no seu caráter seletivo, seja de defesa de projetos elitistas x classe trabalhadora, seja no modo como o estado criminaliza e cria mecanismos de desqualificação das lutas populares.
Alguns fatores parecem concorrer para esse tipo de resultado, a saber : o fato de que a segunda parte foi tratada exaustiva e restritamente a partir de variados exemplos de repressão da ação do estado ; a confusão entre ação criminalizadora do estado e o papel da grande mídia ; o distanciamento do enunciado ao abordar a conjuntura.
Cabe \textbf{lembrar} que a prova se deu em um período pós-eleitoral extremamente conturbado.
Assim, foram inúmeras as referências ao pleito eleitoral, como também foram diversificados os posicionamentos políticos diante daquela conjuntura. - Houve uma ênfase na descrição da importância e da trajetória da democracia em nosso país, apontando o seu surgimento, a partir da Constituição Federal de 1988, desprezando, com isso, a vigência de processos democráticos anteriores à ditadura militar. - Muitas respostas se referiram à ação penal e violenta do estado na vida cotidiana, dirigida aos sujeitos marcados por classe e raça, culminando com críticas à injustiça dos padrões de relação com as classes trabalhadoras e ao sistema carcerário no país, fugindo ao que foi solicitado na questão, possivelmente por uma incompreensão do segundo tópico do comando.
Não raro, os estudantes confundiram"\textbf{criminalização}"com"criminalidade", demonstrando desconhecer a reflexão sobre os processos de \textbf{criminalização} e judicialização da questão social. - O discurso sobre a criminalidade, não raro, resvalou no preconceito contra os movimentos e lutas sociais, veiculando posições favoráveis aos processos de \textbf{criminalização}.
Alguns estudantes demonstraram apreensão acrítica em relação aos movimentos sociais na atualidade ou incorporaram de modo"positivo", em suas respostas, o discurso da \textbf{criminalização} dos movimentos, justificando a sua punição como uma ação necessária do estado, para conter a"desvirtuação"dos objetivos"pacíficos"dos movimentos sociais em"desordem pública".
Nesse caso, as respostas colidiram com o conteúdo previsto no padrão. - Por último, destaca -se o total descolamento dos textos elaborados pelos estudantes em relação ao padrão de resposta, pelo uso recorrente de referências a conteúdos de outras questões da parte objetiva da prova.
Ainda que a intenção dos estudantes, nesse caso, fosse a de se apoiar numa construção textual mais consistente, ou mesmo pertinente ao solicitado na questão, a abordagem elaborada, em quase nada se aproximou dos dois itens requeridos para a resposta da questão.
Cabe salientar que a linguagem utilizada pelos estudantes raramente foi adequada, não apenas pela dificuldade de se expressarem conforme a norma-padrão, mas por usarem jargões e frases feitas, sem significação relevante.
A abordagem mais utilizada para responder à questão foi a que se deteve, de forma correta, no tratamento da importância histórica dos movimentos sociais para a conquista e consolidação dos direitos e da democracia no Brasil.
Essa tendência, relativa à primeira parte da questão, pôde ser observada seja em relação aos textos mais elaborados, com boas digressões históricas, assim como naqueles mais limitados em termos de vocabulário e ideias.
Essa relação acabou se configurando como a solução mais frequente, independentemente das respostas \textbf{estarem} \textbf{certas} ou não.
A abordagem histórica foi a mais frequente para definir a relação entre movimentos sociais e democracia, apontando eventos históricos, especificamente, as lutas para o fim da ditadura civil-militar instalada em 1964 e os movimentos em \textbf{torno} das eleições diretas no período da abertura política no Brasil.
Além desses parâmetros, observaram -se análises das lutas democráticas como expressão da estrutura social desigual, típica do sistema capitalista, numa abordagem sociológica da questão.
De um modo geral, os estudantes apreenderam o significado histórico dos movimentos sociais na construção da democracia brasileira, reconhecendo, em menor ou maior grau de consistência, as ações repressivas por parte do estado.
Alguns conseguiram identificar, com precisão, o caráter seletivo das lutas e da relação estabelecida pelo estado, descrevendo métodos utilizados pelo poder público para controle e cerceamento dos grupos em sua ação política.
A repressão aos movimentos sociais, como parte dos conflitos de classe e a forma histórica com que o estado atuou sobre os mesmos, representou o segundo item do padrão de resposta mais abordado, mas em absoluta distância quantitativa em relação ao anterior.
E ainda, em bem menor grau de frequência, com referência ao caráter “ restrito ” da democracia no Brasil, que foi mais contundentemente abordada quando aludiu à inexistência de uma efetiva democracia no Brasil, em razão do tratamento diferenciado dado às classes sociais e aos cidadãos, de acordo com o seu poder aquisitivo, ou seja, pela cisão entre"ricos"e"pobres", entre os que"têm"e os que"não têm"posses, etc.
No tocante à segunda parte da questão, sobre o caráter seletivo da ação penal do estado na atualidade, o tema mais abordado foi o da repressão às formas de resistência dos trabalhadores, aliado ao das formas de organização dos movimentos pelos diferentes sujeitos coletivos.
Em muitas respostas, os movimentos sociais foram interpretados como dispositivos de pressão popular sobre o estado, tornando -se alvo de dura repressão, tais como o movimento sindical, sem-terra, LGBT, etc.
Os meios de coerção apoiados em ações deslegitimadoras das bandeiras de luta dos movimentos sociais, e o exercício de mecanismos de controle da “ \textbf{ordem} ” foram pouco desenvolvidos nas respostas avaliadas.
Os principais equívocos recaíram sobre o tratamento acrítico dos movimentos, aderindo à noção de violência como justificativa de \textbf{criminalização} da ação política, e, por outro lado, o entendimento da ação penal sobre os sujeitos individuais e não sobre a luta política, como solicitava o enunciado.
Ademais, algumas respostas restringiam a democracia a uma experiência vivida no Brasil pós-Constituição de 1988, ou identificavam, genericamente, todas as lutas sociais como necessariamente direcionadas para a defesa da democracia.
Cabe destacar que uma parte considerável das respostas tratou a relação entre movimentos sociais e democracia como se fossem dois sujeitos ou duas coisas e não de forma a considerar a relação entre sujeito e processo.
Esse equívoco se mostrou, em alguns textos, como uma dificuldade de expressão, e em outros, como uma incapacidade de apreensão histórica.
Chamaram a atenção, também, as respostas em que os movimentos sociais aparecem desprovidos de seu caráter de classe, caracterizados como uma busca pelo bem comum, harmonia social ou, ainda, como avanço da sociedade.
Também se destaca como solução recorrente a abordagem da repressão aos movimentos sociais como “ resposta das autoridades policiais ” aos atos de vandalismo amplamente divulgados pela grande mídia durante as manifestações de rua ocorridas desde 2013.
Houve, aqui, uma enorme confusão entre a compreensão do que trata a questão e da resposta dada pelo estado, já que os estudantes não abordaram a repressão nos termos colocados pelos textos e pelo enunciado, o que levou a respostas extremamente conflitantes, ou seja, afirmando a importância dos movimentos sociais em um primeiro momento e ressaltando que as leis deveriam ser mais duras para com os movimentos que “ transgredissem ” a \textbf{ordem} pública, logo a seguir.
Muitas respostas avançaram nesse item afirmando sobre a necessidade de mudança no aparato legislativo em relação à \textbf{criminalização} dos negros e pobres e a uma enorme variedade de exemplos, clara e amplamente apoiados em temas presentes em outras questões da prova.
Nessa direção, um dos equívocos mais comuns foi o de confundir"\textbf{criminalização}"com"criminalidade", produzindo respostas cujo conteúdo colidia frontalmente com o previsto no padrão de resposta.
Não raro, tal posicionamento veio acompanhado da defesa do endurecimento das penas e do tratamento dado aos prisioneiros, afrontando os direitos humanos e os preceitos do projeto ético e político do Serviço Social.
Outro equívoco importante foi o de vincular o caráter seletivo da ação penal do estado com a questão do sistema penal e do encarceramento, no Brasil, fugindo totalmente ao padrão de resposta.
As referências aos movimentos sociais recaíram sobre o campo das exemplificações : movimentos dos professores, movimentos dos sem teto, dos sem-terra, de mulheres, LGBTs, etc., mas, quase sempre, trazendo para o texto situações vivenciadas nas realidades locais onde estudam ou moram os alunos.
Foi interessante observar que o tema encontrou eco nas diferentes regiões do país, mas ganhou muito mais o ar de denúncia do que de referência que apoiasse a abordagem nos termos requeridos no enunciado.
Foram inúmeras as vezes em que as respostas dadas pelos estudantes não correspondiam ao que \textbf{estava} previsto no padrão de resposta.
Como exemplos, merecem destaque : a idealização da democracia ; a caracterização dos movimentos sociais como expressão da \textbf{criminalização} ; a abordagem da segurança pública ; a \textbf{criminalização} e aprisionamento dos pobres ; a violência promovida pelos movimentos sociais ( e não pelo estado ).
Destaca -se, por último, o fato de que bons textos foram encontrados abordando a dinâmica do modo de produção capitalista, a história da constituição das classes sociais no Brasil e no mundo e as funções coercitivas e educativas do estado, apoiados em autores que compõem a literatura utilizada na formação profissional em Serviço Social, mas que se restringiram a discussões secundárias em relação ao que era solicitado pela questão.
Em termos gerais, o desempenho dos estudantes foi frágil, pois 15 \% dos estudantes deixaram a resposta em branco, e mais da metade das respostas dadas foram consideradas fracas ( 56 \% do total de participantes ).
Essas respostas não abordavam os conteúdos solicitados no padrão de resposta ou, simplesmente, reproduziam trechos do enunciado ou de outras questões da prova.
Muitas vezes, apresentavam uma defesa idealizada da democracia, sem explicitar o seu sentido e sua relação com os movimentos sociais na trajetória histórica da sociedade brasileira.
Houve, ainda, interpretações equivocadas como a confusão com temáticas relativas à criminalidade e à segurança pública e os que tratavam a \textbf{criminalização} do ponto de vista individual, como um comportamento ocasionado pela vigência de uma desigualdade social que, inevitavelmente, o levaria à prática de crimes.
Houve uma expressiva incidência de respostas nas quais a abordagem dos movimentos sociais se fez presente de forma adequada, ainda que nenhum outro item considerado no padrão de resposta fosse tratado de forma correta.
Portanto, as notas mais baixas ( de zero a 30 pontos ) representaram conteúdos inconsistentes, plagiados ou abordaram, basicamente, a ideia de que movimentos sociais são importantes para a construção da democracia.
As respostas consideradas medianas, com notas entre 35 e 70 pontos, foram atribuídas a 26 \% dos participantes.
O mais comum foi enfatizarem o tratamento da importância dos movimentos sociais.
Em relação à segunda parte do comando, em decorrência do excesso de exemplificações e distorções na compreensão do caráter seletivo da ação penal do estado, pouco abordaram os mecanismos de coerção e controle da \textbf{ordem}.
Nas respostas medianas, em geral, os estudantes demonstraram conhecimento de que, historicamente, a luta pela democracia implicou respostas repressivas e controladoras por parte do estado, exemplificando esse processo no passado ou mesmo com indicativos sobre esse padrão de relação no presente.
As respostas com notas consideradas boas ( de 75 a 100 pontos ) conseguiram cobrir a maior parte dos itens que compunham o padrão de resposta, mas foram 3 \% do total de participantes.
Nesses casos, deve ser destacado que o caráter seletivo da ação do estado teve maior ênfase dentre as respostas \textbf{certas} do que os modos de controle da \textbf{ordem} e os meios de coerção que atuam na deslegitimação das bandeiras de luta dos movimentos sociais.
As respostas consideradas boas se diferenciaram pela capacidade de identificar, além do significado histórico dos movimentos sociais na democracia, o caráter seletivo das lutas e das respostas repressivas do estado, indicando exemplos de situações no passado e no presente, e/ou explicitando os mecanismos coercitivos adotados pelo estado, como prisões, ações violentas da polícia, desqualificação dos movimentos e de suas pautas de luta junto à mídia.
Foi marcante a dificuldade de expressão dos conteúdos e conhecimentos por parte dos estudantes, no que tange ao domínio da língua portuguesa.
Muito embora o domínio da língua não se constitua o foco dessa parte da avaliação, é impossível desenvolver algum argumento crítico ou autoral sem o mesmo.
Trata -se de uma preocupação latente em todos os cursos de graduação da área, pois há um visível déficit nos processos de leitura, apreensão e escrita acadêmicas.
Relacionando mais detidamente os conteúdos presentes nas Diretrizes Curriculares, em especial, a relação entre estado, Sociedade Civil e as particularidades do modo de produção de capitalista, os estudantes que realizaram a prova parecem gozar de \textbf{certo} grau básico de apropriação.
As respostas que conseguiram desenvolver algum raciocínio argumentativo expressavam o enquadramento do debate no marco estrutural e contraditório da sociedade capitalista, numa apreensão dos movimentos sociais como sujeito político desse processo.
Cabe destacar, contudo, que, do ponto de vista histórico, essa apreensão se mostrou problemática.
Repetiu -se nas respostas um erro comum nas avaliações dos cursos de graduação em Serviço Social - demarcar o processo de aprovação da Constituição Federal de 1988 como início da democracia, dos direitos sociais e, até, da relação entre o estado e a Sociedade Civil no Brasil.
Não foram poucas as respostas que apresentaram essa compreensão distorcida do significado daquela legislação no percurso da proteção social, dos direitos e da democracia no país.
Além, é óbvio, do significado dos próprios movimentos sociais, tomados como consequência da aprovação da Carta Magna, em muitas respostas.
As respostas dadas pelos estudantes eram pouco consistentes do ponto de vista do domínio histórico, restringindo -se aos marcos relativamente recentes, resvalando em equívocos que, de um lado, demonstram desconhecimento de processos históricos passíveis de comparação em períodos anteriores ; por outro lado, parecem indicar que eventos contemporâneos não têm sido objeto de análise nos processos de formação, pois, algumas vezes, foram tratados na ótica das interpretações midiáticas e do senso comum.
No que tange ao menor rendimento na segunda parte da questão, observou -se que a questão exigia níveis distintos de mediações teóricas, históricas e com a realidade atual.
Ao abordar elementos da realidade atual, os estudantes, de forma geral, perdiam a condição de estabelecer mediações com e a partir das categorias teóricas dos textos, centrando -se mais no reconhecimento de situações concretas, mas amparadas apenas pela empiria.
Certamente, trata -se de uma dificuldade do campo da formação profissional e não exclusivamente desta questão da prova. 6.3.2.3 - Análise da Questão Discursiva 4 do Componente de Conhecimento Específico Na Tabela 6.19, constam as informações relativas à questão 4 do conjunto de questões do Componente de Conhecimento Específico.
O desempenho dos estudantes nessa questão foi inferior ao das questões 3 e 5.
A Média geral do Brasil foi 12,6, sendo a maior Média registrada na região Nordeste ( 17,2 ), e a menor, na região Sul ( 8,8 ).
A nota Máxima ( 100,0 ) foi atingida por, pelo menos, um concluinte de todas as regiões.
A Mediana em todo o Brasil foi zero, o mesmo valor obtido nas cinco Grandes Regiões, indicando que pelo menos 50 \% dos participantes tiraram nota zero.
Também foi zero a nota mínima de todas as regiões.
No Gráfico 6.12, \textbf{está} representada a distribuição de notas da questão discursiva 4, do Componente de Conhecimento Específico.
Essa distribuição tem moda na classe de estudantes que, tendo respondido à questão, obtiveram nota zero, seguido de perto por aqueles que deixaram a questão em branco.
A soma das frequências desses dois intervalos supera 60 \% dos respondentes.
O intervalo ( 10 ; 20 ] caracteriza -se como um máximo local, com cerca de 10 \% do total. 6.3.2.4 Comentários sobre as respostas à Questão Discursiva 4 A Questão 4 articula, principalmente, dois núcleos de fundamentação da formação profissional, conforme as Diretrizes Curriculares.
O núcleo de fundamentos teórico-metodológicos da vida social, que envolve conhecimentos sobre o ser social, particularmente enfocando as mudanças no capitalismo e, na particularidade do estado, com o neoliberalismo e seus desdobramentos sobre a questão social e as políticas sociais.
E o núcleo de fundamentos do trabalho profissional, especificamente, o impacto da contrarreforma do estado sobre o trabalho do assistente social, no quadro das transformações gerais no mundo do trabalho e suas expressões especificamente no âmbito de atuação do Serviço Social.
A questão se articulava com as exigências das Diretrizes Curriculares para os cursos de Serviço Social, ao exigir do aluno a análise histórico-concreta dos processos sociais, sob a perspectiva da totalidade.
Pressupunha, para tanto, o trato teórico e histórico da realidade social, associado à crítica ao processo histórico como totalidade concreta, síntese de múltiplas determinações.
Os conteúdos exigidos transitavam pela articulação entre as esferas da produção e reprodução social.
A análise exigia a compreensão do conjunto de transformações sociais ocorridas mais recentemente no capitalismo, nas quais a dimensão econômica tem centralidade.
Dessa forma, considerar essa dimensão era essencial para explicar as alterações do estado e das classes sociais, bem como as relações entre ambos.
Para redigir uma resposta, na temática da contrarreforma do estado, a compreensão da mais recente crise estrutural do capital importava muito.
Era a partir dela que se poderiam compreender os efeitos sobre a questão social, deflagrados pelas respostas ao cenário de crise econômica, bem como sua incidência na esfera da produção, nos processos de trabalho e, sobretudo, na regulação social, na qual o estado assume papel central.
As transformações na esfera da produção e da reprodução social têm claras incidências na questão social e, portanto, ocupavam espaço relevante no conteúdo da resposta.
Acrescenta -se, ainda, que a Questão 4 exigia considerar a particularidade histórica da sociedade brasileira, ao apontar para a contrarreforma do estado e seus múltiplos efeitos na institucionalidade relacionada aos direitos sociais e, portanto, às políticas sociais.
Da mesma forma, o domínio dos efeitos sobre o mundo do trabalho tornou -se uma exigência, à medida que o trabalho se constitui como uma dimensão fundamental para compreender os nexos e as relações em \textbf{torno} da questão social.
Por fim, a questão apontava, ainda, a repercussão desse quadro sobre o trabalho profissional do assistente social e demandava conhecimentos acerca das mudanças no trabalho, no contexto da reestruturação produtiva e sobre a interrupção de um padrão de regulação social, sob os influxos da flexibilização e da desregulamentação voltados aos interesses da financeirização da economia.
Na Questão 4, pretendia -se avaliar, portanto, um conjunto de conhecimentos e habilidades estabelecidos nas Diretrizes Curriculares, que se referem diretamente ao campo do trabalho como parte da totalidade social, tendo como problemática central a inflexão nas políticas sociais e direitos sociais a partir das prerrogativas neoliberais.
A partir desse quadro, o estudante deveria ser capaz de mencionar inflexões sobre o trabalho do assistente social, em termos do adensamento da questão social com a diminuição das políticas sociais, o aprofundamento do desemprego estrutural e a precarização do trabalho.
A questão requeria que o estudante tivesse capacidade teórica para compreender a contrarreforma do estado, no quadro histórico concreto do neoliberalismo, e competência analítica para evidenciar os impactos das mudanças do capitalismo sobre a profissão, evidenciando, portanto, o trabalho profissional no contexto das relações sociais de seu tempo, seja enfocando a dinâmica das demandas sociais à profissão, seja identificando o modo como esses condicionantes interpelam as condições e as relações de trabalho em que os assistentes sociais \textbf{estão} inseridos.
A questão pode ser considerada de nível médio, tendo em vista que seu enunciado pressupunha a apropriação de conhecimentos sobre estado e transformações contemporâneas, além de uma visão histórica da profissão, incluindo a problematização sobre os limites impostos pelo assalariamento.
A temática levantada, decerto, é de fundamental importância para ser avaliada junto a formandos dos cursos da área.
Muito embora o enunciado da questão tenha sido apresentado de modo conciso, com informação suficiente e adequada para o entendimento da questão formulada, foi possível observar que ele permite uma interpretação dúbia sobre os “ rebatimentos na profissão ”, pois não deixou clara a intenção de que o estudante realizasse considerações sobre os impactos da contrarreforma na profissão, na medida em que a expressão “ rebatimento ” diz respeito a rebater, replicar e conter.
Isso, possivelmente, levou a maioria dos estudantes a abordar como a profissão deveria responder ao neoliberalismo, ou seja, com foco na ideia de responder por meio do projeto profissional, mantendo -se na defesa dos direitos sociais, articulando -se aos movimentos sociais ou aprimorando buscas nas redes de serviços por maiores recursos.
Isso destoa da expectativa de analisar a profissão como condicionada pelas relações sociais em que \textbf{está} inserida e, por isso, tensionada ( ou limitada ) pela contração das políticas sociais no estado, na fase neoliberal.
O padrão de resposta \textbf{estava} bem elaborado e abrangia o conhecimento mínimo esperado pelos estudantes sobre esse tema.
Exigia, por parte do estudante, a perspectiva de análise histórico-concreta, capaz de articular as múltiplas determinações que envolvem o fenômeno analisado, sobretudo, ao exigir do aluno a identificação dos efeitos gerados pela contrarreforma sobre a questão social, portanto, sobre o trabalho e, em particular, o trabalho do assistente social.
Todavia, foi possível observar inúmeras interpretações diferentes ( e também equivocadas ) da que foi prevista no padrão de resposta, como será visto no tópico seguinte.
Isso ocorreu, notadamente na segunda parte da questão, que solicitava dois rebatimentos da contrarreforma do estado no trabalho profissional do Serviço Social.
Numa análise geral, foi possível observar que foram pouquíssimas as respostas que conseguiram apreender o conjunto dos conhecimentos exigidos.
Houve uma clara dificuldade de responder satisfatoriamente ao conjunto da questão.
A primeira parte da questão foi mais explorada pelos estudantes do que a segunda parte, embora tenha havido um percentual grande de respostas incorretas.
Na primeira parte, em que o estudante deveria explicar a contrarreforma e suas implicações, observaram -se as seguintes tendências : 1 ) priorização de aspectos referentes aos efeitos da contrarreforma do estado brasileiro, no que diz respeito à perda dos direitos sociais, ao desmonte das políticas sociais e ao agravamento das expressões da questão social ; 2 ) indicação frequente de que a contrarreforma e/ou o neoliberalismo são as causas do aumento das desigualdades sociais e da pobreza ; 3 ) uma menção relativamente frequente à doutrina neoliberal, suas características e o modo como o neoliberalismo se estabelece no Brasil, a partir dos anos 1990, contrapondo -se à ideia de universalidade e expansão das políticas sociais, a partir da Constituição de 1988 ( bastante mencionada ) ; 4 ) pouca frequência do conteúdo relativo à concepção de estado mínimo e suas características.
Foi possível identificar, também, um conjunto de respostas em que a ideia da contrarreforma foi confundida com a da reforma trabalhista e da reforma da previdência propostas durante o governo Temer.
Nessa direção, boa parte dos estudantes se restringiram a analisar a contrarreforma do estado como um movimento atual, numa visão restrita desse processo, articulando -o ao contexto atual das chamadas"reformas"e, também, à dinâmica recente de desfinanciamento da política de assistência social e da saúde, desde o governo Temer até hoje.
Isso ocorreu sem que os estudantes sublinhassem que, embora essas sejam medidas instituídas no contexto neoliberal, o neoliberalismo remonta à década de 1990 no país e é mais abrangente que as duas reformas frequentemente mencionadas.
Soma -se a essas observações o fato de que algumas respostas atribuíram um caráter positivo à contrarreforma, argumentando ser uma medida que visava ao combate à corrupção e/ou necessária para que o país retomasse o “ caminho \textbf{certo} ”.
Houve, ainda, com alguma frequência, o entendimento equivocado sobre o conceito de contrarreforma, compreendendo -o como um movimento de garantia de direitos para os trabalhadores, em consonância com as formulações constitucionais de 1988.
Tais respostas revelaram total desconhecimento da temática em relevo.
Quanto à segunda parte da questão, foram poucas as respostas que contemplaram os efeitos da contrarreforma sobre o trabalho do assistente social, numa análise que, muitas vezes, não se remetia aos processos sociais relacionados à contrarreforma do estado.
Foram observadas respostas que apontaram quais deveriam ser as providências dos assistentes sociais em face das consequências negativas para a população frente à adoção dos preceitos da contrarreforma.
Mencionaram, por exemplo, que o Serviço Social deveria embasar sua intervenção no Código de Ética Profissional ou no projeto ético e político da profissão, ou aliar -se aos movimentos sociais para lutar contra as forças conservadoras, etc.
Foi perceptível uma gama de respostas muito mais prescritivas de ações/atitudes diante da contrarreforma do que uma análise propriamente dita das repercussões desse processo sobre a intervenção profissional.
Além disso, em algumas respostas, ocorreu a reprodução de enunciados de outras questões e, até mesmo, das alternativas de questões objetivas da prova.
Desse modo, a segunda parte da questão foi mais negligenciada e recebeu abordagens muito aquém das expectativas, concentrando -se, a maioria das vezes, no viés político e ideológico dos profissionais como sujeitos redentores, militantes de resistência ao neoliberalismo ou na via técnica abstrata em que cabe ao profissional levantar dados e realizar procedimentos, a despeito da realidade apresentada como pano de fundo.
Ou seja, verificou -se uma forte tendência a narrativas marcadas pelo messianismo voluntarioso dos profissionais que, com sua ideologia ou com sua técnica, superam as determinações sociais do contexto do trabalho.
A maioria das respostas, portanto, não atendeu ao esperado pelo padrão de resposta.
Poucos estudantes conseguiram abordar o conjunto das dimensões de análise exigido.
Como possíveis causas, são mais prováveis : a incompreensão do enunciado e o desconhecimento sobre o eixo central da pergunta, que, muitas vezes, apareceram de forma associada.
No tocante à incompreensão do enunciado, verificou -se que muitos estudantes tiveram dificuldade em entender o que é a contrarreforma do estado, de modo que alguns evidenciaram positividades desse fenômeno para os trabalhadores.
Outros estudantes abordaram a contrarreforma como se fosse, talvez, a reconceituação do Serviço Social, reforçando uma reflexão muito endógena à profissão e totalmente equivocada, frente ao enunciado da questão.
Isso, demarca um dos principais contrastes com o padrão de resposta, que explicitamente esperava a evidenciação da contrarreforma como promotora de mudanças regressivas no estado.
O desconhecimento da temática ficou evidenciado pelo número de respostas em branco ou com conteúdo irrisório.
Soma -se a isso o fato de que, em parte expressiva das respostas, a repetição do enunciado da própria questão e de outras questões da prova, confirmam, possivelmente, o desconhecimento do tema e a tentativa infrutífera de reunir subsídios para responder à questão, ainda que de modo precário e incerto.
A linguagem dos estudantes era academicamente muito limitada.
Verificou -se forte ênfase de linguagem coloquial, demonstrando dificuldades elementares no uso da norma culta da língua portuguesa, além de frágeis recursos de generalização conceitual ou argumentativa.
Ficou patente, portanto, a dificuldade dos estudantes no uso dos termos específicos, muitas vezes com o predomínio de uma linguagem coloquial e de senso-comum, não condizente com uma dissertação acadêmica, apoiada em análise teórica, sob a influência de referências bibliográficas de obras e autores especialistas no assunto.
Muitas respostas desconheciam o tema do estado neoliberal e as exigências de privatização das empresas públicas, de externalização das atividades assistenciais, de diminuição dos direitos sociais e de ampliação do desemprego.
Na maioria das respostas eram desconhecidos tanto o termo contrarreforma, quanto o tempo histórico do capital em crise e o neoliberalismo.
Em geral, as respostas apresentaram sérios limites de conhecimentos teóricos, conceituais e históricos.
Desconhecimento esse, responsável, talvez, pela narrativa impressionista, contradizendo as diretrizes curriculares que indicam o desenvolvimento de competências para analisar o movimento histórico da sociedade capitalista, particularmente do Brasil, para compreender o estado como parte da dinâmica de requisições econômicas em \textbf{cada} época.
A maioria das respostas demonstrou despreparo dos estudantes em compreender as políticas sociais expandidas ou retrocedidas dentro desse contexto.
As dificuldades demonstradas evidenciam uma acentuada fragilidade formativa para a profissão.
Os que respondiam sobre os rebatimentos desses processos sobre o trabalho profissional, apresentaram supostas"estratégias"de como o profissional deveria atuar diante desse quadro, em vez de fazer referência às consequências, sobre o trabalho profissional, do ajuste estrutural no capitalismo e da redução dos investimentos estatais no campo das políticas sociais.
Desse modo, limitaram a abordagem da profissão, exigida na questão, à afirmação de que o estado neoliberal e os efeitos deletérios sobre os direitos sociais podem ser contrarrestados, por meio dos compromissos do “ projeto ético-político da profissão ”, o que é uma perspectiva a histórica e que destoa do padrão de resposta, na medida em que desgarra a profissão das determinações sociais.
Portanto, ficou evidente a dificuldade dos estudantes demonstrarem conhecimento de processos sociais mais amplos nos quais a profissão se insere.
Por isso, muitos tentaram explicar o contexto da contrarreforma do estado por meio de narrativas ideopolíticas e aparatos técnicos próprios do Serviço Social, o que, sem dúvida, é uma distorção perante as diretrizes curriculares.
O desempenho dos estudantes na Questão 4 foi bastante frágil.
A média foi 1,3 e foi elevado o número de respostas que não atenderam ao previsto no padrão de respostas.
A correção revelou um significativo desconhecimento dos estudantes sobre a questão formulada e quanto aos conceitos nela imbricados, o que é motivo de preocupação, pois o tema da contrarreforma do estado brasileiro e seus impactos sobre as políticas sociais e o trabalho profissional são conteúdos de alta relevância para a formação do assistente social.
Quase 30 \% dos participantes deixaram a resposta em branco.
Dentre as respostas válidas corrigidas, quase 80 \% foram fracas ( 58 \% do total de participantes ), com notas entre zero e 30 pontos.
As respostas da Questão 4, em sua maioria, \textbf{estão} abaixo da qualificação razoável, demonstrando um desempenho aquém das competências exigidas pelas diretrizes curriculares e pelo edital do Enade-2018.
As respostas medianas ( entre 35 e 70 pontos ), em grande parte ( em \textbf{torno} de 13 \% do total de participantes ), foram respostas incompletas.
Uma parcela significativa elaborou, apenas, um texto sobre a contrarreforma do estado brasileiro, sem mencionar a segunda parte da questão.
Outra parcela buscou analisar as repercussões da contrarreforma no trabalho do assistente social, porém, como sinalizado anteriormente, predominaram respostas prescritivas sobre como, na atualidade, o assistente social deve atuar no âmbito das políticas sociais, diante da perda de direitos causada pela contrarreforma.
As respostas com maior pontuação ( entre 75 e 100 pontos ), apesar da baixa frequência ( menos de 2 \% dos participantes ), abrangeram, com maior ou menor profundidade, os conteúdos previstos no padrão de resposta.
Nessas respostas, os estudantes conseguiram tangenciar o processo de contrarreforma do estado e suas implicações, sob os auspícios da doutrina neoliberal, além de identificar os anos 1990 como marco temporal de um processo contrário à perspectiva da seguridade social e da universalização dos direitos sociais constantes da Constituição de 1988.
Nessas raras respostas, um número reduzido de estudantes teve capacidade para identificar e citar, com base em análise teórica apropriada, os impactos da contrarreforma no trabalho profissional do assistente social.
Em termos gerais, os estudantes demonstraram deficiências nas competências e habilidades exigidas pelas diretrizes curriculares, em especial, nos conhecimentos teóricos e históricos sobre as transformações do capitalismo e as inflexões no estado, com a consequente desregulamentação do trabalho e dos direitos sociais e seus desdobramentos severos sobre as políticas sociais.
Surpreendeu consideravelmente que os estudantes apresentassem grande dificuldade para tratar de um tema bastante difundido nos espaços acadêmicos.
A maioria demonstrou que não sabe o que é a contrarreforma e/ou o neoliberalismo e, consequentemente, desconhecem suas expressões nas políticas públicas.
Além disso, as respostas mostraram uma visão messiânica de que o assistente social deve ser o garantidor dos direitos sociais, tendo como referência, apenas, o exercício do profissional e o código de ética profissional.
O debate acerca do projeto ético político aparece nas respostas como substituto da argumentação teórica e metodológica necessária para o enfrentamento dos desafios profissionais.
Ao que tudo indica, esse equívoco \textbf{está} relacionado a uma dada opção pela solução da questão e não a uma possível falta de clareza do enunciado.
Isso pode indicar que os cursos não \textbf{estão} conseguindo vertebrar o ensino em \textbf{torno} da compreensão sobre a estrutura da sociedade capitalista e da desigualdade social inerente a ela, nos diferentes contextos.
Essa limitação ganha maior consequência por ser um campo 6.3.2.5 Análise da Questão Discursiva 5 do Componente de Conhecimento Específico Na Tabela 6.20, constam as informações relativas à questão 5 do conjunto do Componente de Conhecimento Específico.
O desempenho dos estudantes de todo o Brasil nesta questão foi o melhor entre as três questões discursivas desse componente.
A nota Média dos estudantes de todo o Brasil foi 27,4.
A maior Média foi registrada na região Nordeste ( 30,3 ), enquanto a menor Média foi registrada na região Sul ( 25,2 ).
Quanto à variabilidade das notas, o Desvio padrão dos alunos do Brasil, como um todo, foi 25,7.
Enquanto o maior desvio foi encontrado na região Nordeste ( 27,0 ), o menor foi encontrado na região Sul ( 23,6 ).
A Mediana para o Brasil e na maioria das regiões foi 20,0 ; com exceção das regiões Norte e Nordeste, onde foi 30,0.
Para o conjunto de alunos de Serviço Social do Brasil, a nota Máxima foi 100,0.
A nota Mínima do Brasil foi zero, o mesmo valor encontrado em todas as regiões. de conteúdos estruturador da profissão, capaz de elucidar o sentido social da mesma e as próprias características das demandas profissionais na realidade contemporânea.
A incompreensão desse núcleo temático expõe sérios limites ao exercício profissional competente e socialmente referenciado como almeja o projeto acadêmico na área.
Portanto, as respostas revelam um alerta fundamental para o debate avaliativo da execução dos projetos pedagógicos nas diferentes unidades de ensino.
As respostas analisadas também demonstram que o Ensino Superior não \textbf{está} sendo capaz de superar as limitações formativas dos estudantes, que chegam a esse curso superior com muitas lacunas, incluindo debilidades de leitura e compreensão de textos, assim como na elaboração de narrativas condizentes com uma prova de conhecimentos acadêmicos.
Fica expressa a imaturidade para fazer provas, o que é significativo, na medida em que são estudantes em conclusão de curso.
Outrossim, não é demais sublinhar que as atividades profissionais exigem variadas práticas de textos argumentativos, que requerem consciência sobre causa e efeito, e uma satisfatória articulação linguística, para extrair súmulas conclusivas.
Essas inconsistências da formação geral e do Ensino Superior dos formandos poderão repercutir negativamente na qualidade do trabalho que vierem a desempenhar, ao mesmo tempo que expressam a própria fragilidade do direito à educação no país. 6.3.2.5 Análise da Questão Discursiva 5 do Componente de Conhecimento Específico Na Tabela 6.20, constam as informações relativas à questão 5 do conjunto do Componente de Conhecimento Específico.
O desempenho dos estudantes de todo o Brasil nesta questão foi o melhor entre as três questões discursivas desse componente.
A nota Média dos estudantes de todo o Brasil foi 27,4.
A maior Média foi registrada na região Nordeste ( 30,3 ), enquanto a menor Média foi registrada na região Sul ( 25,2 ).
Quanto à variabilidade das notas, o Desvio padrão dos alunos do Brasil, como um todo, foi 25,7.
Enquanto o maior desvio foi encontrado na região Nordeste ( 27,0 ), o menor foi encontrado na região Sul ( 23,6 ).
A Mediana para o Brasil e na maioria das regiões foi 20,0 ; com exceção das regiões Norte e Nordeste, onde foi 30,0.
Para o conjunto de alunos de Serviço Social do Brasil, a nota Máxima foi 100,0.
A nota Mínima do Brasil foi zero, o mesmo valor encontrado em todas as regiões.
No Gráfico 6.13, é representada a distribuição das notas da questão discursiva 5 do Componente de Conhecimento Específico.
Essa distribuição, como todas as relativas às questões discursivas dessa parte da prova, também tem modal o intervalo correspondente aos que deixaram a questão em branco, com 18 \%, do total.
Somando -se a frequência desse intervalo com a daqueles que obtiveram nota zero chega -se próximo a 36 \%, dos participantes. 6.3.2.6 Comentários sobre as respostas à Questão Discursiva 5 A Questão 5 enuncia um tema emergente no Serviço Social, presente em uma campanha nacional do Conselho Federal de Serviço Social, em conjunto com os Conselhos Regionais de Serviço Social, intitulada “ Assistentes Sociais no combate ao racismo ”, cuja abordagem vem ganhando relevância na agenda político-profissional, nos eventos e publicações das entidades representativas da categoria.
O debate sobre os estudos das relações étnico-raciais no Serviço Social é recente e ainda não possui centralidade nas grandes discussões fomentadas pela área de conhecimento e pela categoria profissional.
Entretanto, ultimamente, esse debate vem tomando corpo através da produção de conhecimento na área – nos eventos acadêmicos e científicos, nas publicações especializadas, além da inclusão de seu conteúdo no ementário de disciplinas curriculares.
Muito embora haja a inclusão desse debate na formação profissional, não são todas as unidades de ensino que possuem docentes especializados e capacitados para o trato dessa questão.
Por isso, cabe sinalizar a importância da Questão 5 ter abordado essa temática, em consonância com o avanço do debate sobre o racismo e as lutas antirracistas no Serviço Social.
A questão, em relação à sua formulação, vincula -se aos conteúdos relacionados às competências e habilidades requeridas aos assistentes sociais, de acordo com as Diretrizes Curriculares do Serviço Social.
Ela o faz a partir de duas dimensões : uma relacionada ao conjunto de determinantes e condicionantes da formação social brasileira, retomando sua constituição histórica, especialmente reportada ao escravismo, mas, também, em sua atualidade.
A outra se relaciona às requisições postas ao trabalho profissional nas lutas de combate ao racismo, de acordo com o projeto ético e político do Serviço Social.
Em relação à constituição histórica da formação social brasileira, era essencial que os estudantes se reportassem ao escravismo como elemento ineliminável de seu passado histórico.
Já em relação à atualidade, era necessário se reportar a como esses traços constitutivos da formação brasileira se reatualizam e se manifestam frente à desigualdade social, que se especifica na questão da raça e/ou da etnia.
Isto é, como se expressam hoje, em termos da exploração e dominação das classes subalternas, tomando -se o recorte da questão racial como, por exemplo : por meio da subalternização e/ou pela desigualdade racial, seja na questão salarial ou de renda ; na violência, na vitimização e/ou aprisionamento que acometem mais diretamente os estratos populacionais, marcados pela desigualdade social, em sua especificidade relacionada à raça e/ou à etnia.
Esses traços podem ser explorados, também, em relação ao mundo do trabalho, em suas particularidades nacionais, tais como nos desníveis salariais e de renda, nas condições e dificuldades de inserção no mercado de trabalho e/ou, propriamente, no desemprego.
Também, suas expressões se vinculam ao acesso à educação, especialmente no que se refere à sua reduzida efetividade de acesso ao Ensino Superior.
E, ainda, no que se refere à baixa representatividade política, ao desrespeito à cultura, à religião e à herança histórica e em seus vínculos também marcados pelas desigualdades de gênero, que expressam maior vulnerabilidade das mulheres negras, entre outras refrações e/ou expressões da “ questão social ”, na formação social brasileira, considerando -se a variável da questão racial e ou étnica.
A segunda dimensão da questão se voltava para a articulação desses conteúdos às requisições feitas aos assistentes sociais, de acordo com o projeto ético e político do Serviço Social.
Isto é, o estudante deveria demonstrar o papel dos assistentes sociais quando, em seu trabalho, vincula -se à luta contra o racismo, como por exemplo : “ a defesa intransigente dos direitos humanos e/ou o combate a todas as formas de preconceito, discriminação e racismo ” ; a “ promoção da igualdade racial e/ou o fortalecimento das políticas, ações afirmativas e iniciativas voltadas à garantia de direitos da população negra ” e/ou em sua “ articulação com os movimentos sociais e/ou defesa da população negra contra intolerâncias, injustiças e violações de direitos ”.
A questão, portanto, era de extrema relevância para a formação profissional.
Uma série de esforços da categoria, tanto no ponto de vista organizativo, quanto na formação profissional, busca a incorporação da dimensão ético-racial na interpretação sócio-histórica da sociedade brasileira, na análise das demandas profissionais e nas possibilidades de respostas institucionais às diversas expressões do racismo.
Analisar a questão social no Brasil, a partir do recorte de gênero, raça/etnia e classe social, seus fundamentos teóricos e repercussões para o trabalho profissional é um objetivo importante das diretrizes curriculares do Serviço Social.
A questão trazia muitas informações em um único enunciado, o que facultava a referência a temas específicos nela contidos, como por exemplo e especialmente, o vínculo da questão racial, também, à questão de gênero.
Ao abordar as desigualdades econômicas, o enunciado explicitava a relação da discussão do racismo com a questão de gênero e, também, como o racismo se expressa no âmbito jurídico formal, mencionando o racismo social e institucional.
Soma -se a esses aspectos indicados no enunciado o uso de uma imagem que também servia de indicativo para a resposta à questão.
Quanto à utilização da imagem na formulação da questão, é possível inferir que a mensagem subjacente a ela não expressava, de modo substantivo, o conteúdo da questão formulada, seja por se referir, especificamente, ao Dia Nacional da Consciência Negra, seja por explicitar o slogan"com racismo não há liberdade", suscitando a abordagem de outros conteúdos por parte dos estudantes ou, mesmo, o aporte das respostas em enunciados de outras questões da prova, como a questão objetiva de nº 14 que tratava da violência letal e do encarceramento de jovens negros.
Interessante observar que, mesmo o enunciado da questão indicando reflexões sobre a particularidade da condição da mulher negra, a imagem a coloca em segundo plano.
A face do homem negro ocupa um maior espaço no quadro da imagem, remetendo a uma prevalência, em termos imagéticos, da figura masculina.
Além do mais, a imagem não revela, em si mesma, potencialidades suficientes para uma análise do racismo como um elemento estruturador da sociedade brasileira e de suas desigualdades históricas, já que, como peça de divulgação institucional, põe ênfase na campanha do conjunto CFESS/CRESS no combate ao racismo.
Nessa imagem, portanto, os"slogans"se sobrepõem à própria imagem, sem oferecer elementos suficientes para induzir os aspectos exigidos no padrão de resposta.
O enunciado, como um todo, fornecia pistas fundamentais para o desenvolvimento da resposta, porque apresentava dados relativos às desigualdades raciais, à perspectiva de que não há democracia racial no país e sinalizava que o racismo é parte inerente da formação social brasileira.
Portanto, o que foi solicitado ao estudante era claro e preciso.
Os textos que compunham o enunciado denotavam o tipo de análise esperada nas respostas, possibilitando que o padrão de resposta pudesse ser alcançado, pois faziam referência aos conteúdos gerais e particulares relativos às expressões do racismo na formação social brasileira e ao trabalho do assistente social.
O padrão de respostas indicou alguns aspectos interessantes para análise de tendências presentes hoje, na profissão, principalmente sob a ótica da formação profissional.
A primeira delas se referia à existência de um projeto ético e político articulado ao trabalho do assistente social.
Mesmo que não se tenha sido observado nas respostas, com uma frequência alta, o reconhecimento dessa proposição profissional, observou -se o indicativo de que aqueles que o mencionam já o relacionam às diretrizes profissionais.
Deve ser considerado, entretanto, que a questão racial não é suficientemente ou adequadamente aprofundada em boa parte das unidades acadêmicas.
Avalia -se que não há, em todas as unidades de ensino, disciplinas obrigatórias ou optativas sobre o tema, ou seja, esse conteúdo é ausente ou incorporado precariamente aos currículos dos cursos de graduação em Serviço Social.
Esse dado leva a considerar que o nível de exigência da questão era, no mínimo, médio.
O grau de dificuldade se associava, sobretudo, aos conhecimentos requeridos e à capacidade de síntese exigida do estudante, tendo em vista espaço e tempo reduzidos para responder à questão no contexto da prova do Enade, com um total de 40 questões, sendo cinco discursivas.
O estudante teria que discorrer sobre o racismo e sua natureza histórica, articulando -o com os dilemas étnico-raciais contemporâneos que vêm sendo muito publicizados na literatura, nos movimentos sociais e em diferentes mídias.
Além disso, deveria, também, estabelecer a correlação desse conteúdo com as propostas de atuação profissional para o seu enfrentamento, o que exigia uma reflexão mais \textbf{apurada} desenvolvida na formação profissional.
O grau de dificuldade também \textbf{estava} associado ao fato de que o estudante deveria articular o seu entendimento sobre racismo na sociedade brasileira ao modo como o assistente social enfrenta essa questão em seu trabalho cotidiano, considerando as bases éticas e políticas que direcionam o exercício profissional.
Ocorre que esse entendimento, muitas vezes, é atravessado pelo racismo arraigado no pensamento conservador e discriminatório, às vezes expresso pelo próprio estudante, como foi identificado em algumas respostas.
Isso é motivo de preocupação, pois uma postura velada ou francamente racista por parte de estudante concluinte do Curso de Graduação em Serviço Social colide frontalmente com os preceitos do Código de Ética Profissional.
Observou -se que os estudantes, de uma forma geral, não demonstraram possuir a gama de conhecimentos necessários para responder à questão formulada.
Quando houve menção aos conceitos previstos, eles se revelaram parcialmente compreendidos.
Um conjunto significativo das respostas apresentou limitações de entendimento do enunciado.
Muitas expressavam um conteúdo idealizado, prescritivo, sustentado no senso comum.
O aspecto jurídico-legal \textbf{esteve} muito presente nas alternativas para a luta contra o racismo, mas indicando, de forma superficial, a sua construção e efetivação no trabalho profissional.
Nas respostas incompletas, no que tange à primeira parte da questão, referente à abordagem do racismo na formação social brasileira, observaram -se as tendências abaixo relacionadas. - Parte significativa das respostas não considerou o racismo no Brasil sob uma perspectiva crítica e histórica.
O escravismo nem sempre foi mencionado e sua \textbf{localização} na atualidade dava a entender que esse fenômeno surgiu desvinculado dos aspectos sócio-históricos e culturais.
Muitas vezes, os estudantes se reportaram, genericamente, ao colonialismo e ao “ passado histórico ”, mas sem especificação do processo de escravidão ou à exploração da força de trabalho da população negra por meio do escravismo, como previsto no padrão de resposta.
Também \textbf{esteve} presente, com menor participação, a referência à escravidão como fenômeno presente até hoje, sem uma melhor demarcação teórico-conceitual e/ou histórica, pois houve equívocos na delimitação do período histórico da escravidão ( por exemplo, remetendo ao “ século passado ”, a “ décadas atrás ”, etc. ).
O escravismo, portanto, foi tomado de forma adjetiva e não substantivamente. - Houve uma grande incidência da vinculação da questão de gênero à questão racial – especialmente considerando a inserção distinta e desigual entre homens e mulheres – e mulheres negras – no mercado de trabalho.
Embora, a rigor, o debate acadêmico, não apenas no Serviço Social, mas também nas Ciências Sociais, trabalhe de forma articulada as questões de classe, raça/etnia e gênero, o que se verificou no conjunto das respostas foi, mais substantivamente, a articulação entre raça e gênero mais vinculada ao enunciado da questão e não propriamente a uma apropriação de conhecimento desse debate. - Foi, também, frequente, relativamente ao conjunto das respostas, o tratamento do racismo como equivalente a outros elementos de diversidade, como a relacionada ao gênero, não só especificamente à maior exploração/dominação das mulheres, mas à diversidade sexual em geral, considerando, também, aspectos de homo/transfobia.
Isto é, um conjunto mais expressivo de respostas articulou distintas diferenças : entre homens e mulheres, entre negros, brancos e indígenas, entre heterossexuais e homossexuais, ressaltando seu alijamento e/ou sua inserção desigual na sociedade. - Houve uma significativa incidência de respostas que citavam situações de racismo na atualidade, porém pautadas no mito da democracia racial.
Muitos estudantes reproduziram esse mito, apontando que só existe uma única raça, “ a humana ”, desconsiderando a luta antirracista organizada pelos movimentos sociais negros presente na sociedade brasileira, com protagonismo na luta pela garantia de direitos e igualdade racial.
Nesse viés, foram expostas críticas às políticas de reparação, como é o caso da lei de cotas raciais, e não só na universidade.
Muitas respostas defendiam que essas políticas afirmativas não deveriam existir, tomando como referência uma suposta igualdade jurídico-formal, uma vez que “ somos todos iguais ” na Constituição Federal, ou por interpretações equivocadas e de fundo moral e religioso sobre o Código de Ética Profissional. - Muitas respostas adquiriram um caráter confessional pautado na vivência pessoal do racismo nas universidades, no mercado de trabalho, na aparência pessoal e nas abordagens policiais.
Não em maior monta, mas também presente, foram as respostas que se reportaram à questão com conteúdo moral e/ou apelando a sentimentos, emoções e/ou à religiosidade ( por exemplo, “ o racismo é uma questão que se aprende em casa ” ; o “ amor e respeito ao próximo ”, “ somos todos filhos de Deus ” ).
O discurso religioso, portanto, também apareceu circundando e/ou tangenciando, de alguma forma, as respostas, ao pôr em relevo o quanto o racismo é incompatível com Deus e com práticas de fraternidade.
Cabe registrar que, de modo residual, em algumas poucas respostas foi possível identificar um tipo de pensamento conservador que procura esvaziar as perspectivas que sinalizam para uma crítica sistêmica ao racismo.
Reduzida meramente a críticas comportamentais, sem qualquer fundamentação teórica e/ou histórica, a narrativa dessas respostas busca direcionar a causa do racismo às próprias vítimas do racismo : mulheres negras e homens negros.
Assim, veiculam a ideia de que a culpa do racismo é do próprio negro.
São crentes no “ racismo reverso ” e na suposta “ democracia racial ”.
São críticos ferrenhos à política de cotas.
São aqueles que sustentam que não existe desigualdade entre raças por força da miscigenação brasileira ou, ainda, que dizem que só existe uma raça, a “ humana ”.
No que se refere à segunda parte da questão, relativa à importância da luta contra o racismo para o trabalho de assistentes sociais, podem ser apontadas, nas respostas, as seguintes tendências abaixo descritas : - Houve uma maior presença de respostas vinculadas apenas, genericamente, à afirmação do necessário combate ao racismo por parte do assistente social, sem maior explicitação dos requerimentos pertinentes ao exercício profissional.
Os estudantes responderam, de forma breve e rasa, repetindo jargões contra o racismo assentados, muitas vezes, em visões humanistas ou em"palavras de \textbf{ordem}", mas com pouco aprofundamento crítico acerca das formas concretas de enfrentamento do racismo no trabalho dos assistentes sociais.
Deve ser aqui destacado que foram pouco mencionadas as outras duas possibilidades previstas pelo padrão de resposta.
De modo incipiente aparecerem respostas que se referiam ao fortalecimento das políticas, a ações afirmativas e/ou a iniciativas voltadas à garantia de direitos ; e mais raras, ainda, foram aquelas que mencionaram a articulação com os movimentos sociais. - Uma significativa parcela de respostas se deteve em relacionar a luta contra o racismo ao projeto ético e político do Serviço Social de modo genérico.
E, quando houve explicitação dos vínculos à dimensão ético-política do trabalho profissional, o maior número de respostas se voltou à afirmação de princípios do Código de Ética Profissional – “ Exercício do Serviço Social sem ser discriminado/a, nem discriminar, por questões de inserção de classe social, gênero, etnia, religião, nacionalidade, orientação sexual, identidade de gênero, idade e condição física ” ; e/ou “ opção por um projeto profissional vinculado ao processo de construção de uma nova \textbf{ordem} societária, sem dominação, exploração de classe, etnia e gênero ”. ( CFESS,1993 ).
Ressalta -se que, nem sempre, esses princípios foram enunciados de forma integral, mas apenas em uma formulação mais geral de seu conteúdo.
Houve menção recorrente ao Código de Ética Profissional e uma frequência significativa de suas orientações éticas e políticas, principalmente àquelas relativas à diversidade étnica, racial, de gênero e aos direitos sociais, políticos e humanos.
Notou -se, no entanto, que as respostas se restringiram aos princípios do Código de Ética e não ampliaram a discussão dos aspectos relacionados ao projeto ético-político como um todo. - Foram identificados, em boa parte das respostas, padrões repetitivos e fragmentados quanto à compreensão do trabalho do assistente social na perspectiva do projeto ético e político do Serviço Social, já que a mediação, enquanto categoria central no trabalho profissional, não ocupou a centralidade nos debates, principalmente aqueles afeitos diretamente à experiência concreta.
Depreende -se que esse aspecto ainda se constitui um desafio à profissão, tendo em vista que não parece \textbf{estar} completamente compreendido entre os estudantes.
O fato de terem associado, de modo recorrente, o Código de Ética ao projeto profissional também sinaliza o reconhecimento da existência do debate na formação profissional, ou seja, indica que os estudantes conseguem estabelecer algum nível de articulação com essas diretrizes, ainda que a mediação do trabalho profissional seja insuficiente.
A debilidade dessa apreensão foi evidenciada, também, por outro equívoco, o de confundir a profissão de Serviço Social com a política pública de Assistência Social, pois, não raro, os estudantes usaram a expressão"profissional de assistência social". - Um outro aspecto foi a presença de uma valoração moral nas respostas, de como, através do trabalho profissional, se poderia enfrentar e combater o racismo.
Há, ainda, no imaginário dos estudantes, alguns ranços conservadores a despeito dos avanços realizados nas últimas décadas no projeto profissional.
Em alguns casos, as respostas foram baseadas no senso comum, com base em um imaginário marcado pela visão moral e escravocrata, que ainda vigora na sociedade brasileira. - E, finalmente, com alguma participação, foram observadas respostas que apenas reafirmaram a importância da atuação do Conselho Federal de Serviço Social ( CFESS ) ao promover o debate sobre a questão racial, destacando o Dia da Consciência Negra, talvez por influência da imagem utilizada, de forma ilustrativa, no enunciado da questão.
Outrossim, a linguagem adotada pelos estudantes, em sua maioria, foi adequada em relação aos termos pertinentes e usuais no Serviço Social.
Todavia, como nas respostas às outras questões discursivas, os textos continham muitos erros de ortografia e gramaticais, bem como de argumentação com coerência e organização textual, o que surpreende do ponto de vista da formação profissional, na medida em que a maioria dos estudantes é concluinte e, em tese, aprendeu a elaborar trabalhos acadêmicos os mais variados, além de documentos técnicos profissionais nos espaços institucionais em que realizaram estágio supervisionado.
Chama atenção a alta recorrência de respostas que parafrasearam ou reproduziram integralmente o enunciado da questão e que, por esse motivo, não avançaram no processo de reflexão sobre o racismo e a importância da luta contra o racismo para o trabalho dos assistentes sociais.
Foi preocupante o quantitativo de respostas com nota zero, apesar de inferior ao das questões 3 e 4, principalmente considerando todos os elementos presentes no enunciado.
Houve, também, uma alta incidência de respostas incompletas, que contemplaram o conteúdo da primeira parte e não abordaram o conteúdo da segunda parte da questão e vice-versa.
Com relação à primeira parte da questão, a maioria dos estudantes que abordou o tema da escravidão fez de forma superficial, situando, nas suas respostas, algum conteúdo sobre o passado histórico da colonização e a presença dos negros nesse cenário.
Notou -se a ausência do debate sobre o racismo como elemento estrutural na formação social brasileira e uma concentração da abordagem sobre as diferentes expressões contemporâneas do racismo, bem como sobre a realidade atual do negro na sociedade brasileira, com destaque para a exclusão do negro no mercado de trabalho e seu encarceramento no sistema prisional ( certamente como influência de uma questão objetiva ).
Houve muita confusão relacionada ao mito da democracia racial.
Os estudantes negaram a sua existência ( conforme consta na citação presente no enunciado da questão ), mas atribuíram a explicação do problema, sobretudo, à dimensão individual, ou seja, quem é racista o é por uma questão pessoal.
É possível evidenciar esse discurso através de exemplos como, o de que todos os brasileiros possuem “ sangue negro ”, que “ todos somos misturados ”, que a “ cor da pele não significa nada ”, que “ ser negro não tem nada a ver com caráter ” e que, por esse motivo, não há necessidade, motivo ou sentido em"sermos racistas".
As respostas sugerem que o Brasil possui todas as condições para não ser um país racista, considerando a miscigenação, mas que apesar dela, o racismo ainda vigora.
Isso mostra a pouca densidade crítica das respostas que não ancoram o escravismo e o racismo na estrutura de desigualdade social no Brasil.
Verificou -se, também, uma confusão entre raça, como um conceito sociológico, e “ cor ”, como um elemento de caracterização racial.
Na segunda parte, ocorreu semelhante processo.
Grande parte das respostas citou, de maneira direta ou indireta, o combate a todas as formas de preconceito, discriminação e racismo como um princípio presente no projeto ético e político da profissão.
Entretanto, a maior dificuldade foi identificada em outras dimensões inerentes ao trabalho profissional, como a promoção de políticas e programas direcionados à garantia de direitos da população negra e a articulação com os movimentos sociais e/ou de defesa dessa população.
As formas de enfrentamento mediante o desenvolvimento de políticas sociais, que incorporem a dimensão racial, foi o elo mais frágil da resposta.
Em geral, nesse segundo item, ou apareceu o princípio ético de combate a todas as formas de preconceito, ou a menção à defesa dos direitos da população negra.
Muitas vezes, as respostas sinalizam que não há o que o Serviço Social possa fazer, porque o racismo é algo inerente a algumas pessoas e, portanto, as ações profissionais são inócuas.
Essa interpretação se refere ao mito da democracia racial, já sinalizado anteriormente.
O percentual de respostas em branco foi de 18 \% do total de participantes.
As respostas na faixa de pontuação entre zero e 30 pontos ( 43 \% do total de participantes ), podem ser divididas entre : aquelas que desenvolveram conteúdos incipientes para as duas partes da questão ; aquelas que desenvolveram conteúdos somente para uma das duas partes ; e, ainda, aquelas que não responderam a nenhuma das partes da questão ou fugiram ao tema.
As respostas que obtiveram pontuação média, na faixa de 35 a 70 pontos, foram as de 36 \% dos participantes.
Em sua maioria, conseguiram abordar o conteúdo das duas partes da questão de forma mais substantiva e consistente.
Os principais acertos se deram quando a resposta \textbf{estava} diretamente relacionada ao enunciado da questão, ainda que alguns conteúdos previstos no padrão de resposta não fossem abordados ou explicitados.
As respostas, mesmo apresentando alguma debilidade teórico-conceitual, conseguiram não só discorrer sobre o racismo no Brasil, citando situações de discriminação, subalternização e de desigualdade, expressas em inúmeras variáveis da vida social e econômica da população negra, como, também, abordaram a importância do combate às variadas formas de preconceito e discriminação, reconhecendo que o projeto ético e político do Serviço Social \textbf{está} associado à luta contra toda e qualquer forma de racismo.
Destacaram -se, no conjunto das respostas, aquelas que conjugaram parcial ou integralmente a primeira parte da resposta ( relações entre racismo e formação social brasileira : no passado e/ou na atualidade ), com algum elemento da segunda parte ( requisições às lutas/combate ao racismo no trabalho profissional do assistente social ) ou vice-versa.
As respostas que obtiveram pontuação elevada, na faixa entre 75 e 100 pontos, foram menos expressivas ( próximo a 3 \% ).
Em sua maioria, as respostas \textbf{estavam} fundamentadas com os conteúdos curriculares do Curso de Graduação e atenderam, de modo substantivo, ao previsto no padrão de resposta.
Apesar de as respostas boas terem figurado em menor número, tais respostas são indicativas de que o processo de renovação ocorrido no Serviço Social brasileiro, tanto no âmbito acadêmico, como na produção de conhecimento e na organização da categoria, tem produzido, ainda que minoritariamente, uma geração de futuros profissionais com elevado grau de competência e compreensão dos desafios profissionais.
O desempenho dos estudantes nessa questão demonstrou, em linhas gerais, que há lacunas no processo formativo, pois foi baixo o quantitativo de estudantes que demonstrou capacidade para identificar, com base em análise teórica apropriada, as características do racismo no Brasil, a compreensão dos seus rebatimentos atuais e os desafios ao trabalho do profissional formado em Serviço Social.
Os estudantes entendem o racismo como objeto de atuação profissional, mas não avançam nessa direção.
Em sua maioria, citam um dos princípios do Código de Ética, mas com reduzido aprofundamento sobre o papel profissional no combate ao racismo e na garantia de direitos da população negra.
O debate acerca do projeto ético e político aparece, muitas vezes, como substituto da argumentação teórica e metodológica necessária para o enfrentamento dos desafios profissionais.
Observou -se que há limitações na apreensão de aspectos importantes, seja no que se refere à perspectiva histórica da formação social brasileira, quanto no que se refere às requisições que são postas ao trabalho dos assistentes sociais, de acordo com o projeto ético e político profissional.
Sendo assim, a correção dessa questão suscitou indagações sobre as condições da formação dos estudantes, tendo em vista que a questão trata de um tema importante para o Serviço Social.
Por isso, é necessário que as unidades de ensino ampliem o debate sobre essa temática nos Cursos de Graduação, pois as respostas dos estudantes indicam a necessidade de maior discussão sobre a dimensão histórica do racismo na formação social brasileira e, particularmente, sobre a luta contra o racismo como uma construção coletiva e não limitada ao âmbito de uma categoria profissional.
Também é necessário aprofundar, na formação profissional, a articulação com os movimentos organizados e trabalhar as representações ainda existentes no imaginário da sociedade ( e dos estudantes ) sobre o racismo, bem como os conteúdos valorativos e religiosos que evocam a presença do religioso, de “ Deus ”, da moral, no trato da questão racial, colidindo com o processo de laicização da profissão.
Cabe, também, ressaltar que a compreensão do projeto profissional ainda não \textbf{está} totalmente absorvida, de modo crítico, entre os estudantes, sendo necessário investir em conteúdos programáticos nas disciplinas dos currículos que propiciem superar jargões e clichês em sua apreensão.
De modo geral, evidenciou -se, nessa correção, que há pouca profundidade teórica, metodológica, técnica e interventiva que sustente as análises críticas do fenômeno e dos mecanismos imprescindíveis para seu equacionamento.
Mais uma vez, é necessário registrar que o tema tratado, apesar de \textbf{estar} presente na agenda profissional, ainda é um assunto pouco explorado ou é abordado, de forma superficial e precária, nas Instituições de ensino.
Pode ser que a fragilidade do desempenho estudantil \textbf{esteja} também relacionada ao fato de que essa temática foi recentemente incorporada como objeto de disciplinas nas grades curriculares dos cursos de graduação em Serviço Social, demandando maior aprofundamento das discussões relativas ao racismo. 6.3.3 Considerações Finais A experiência de correção das questões discursivas do Enade-2018 suscita sérias indagações sobre as condições da formação profissional do assistente social na atualidade, tendo em vista que as temáticas abordadas integram a grade curricular dos Cursos de Graduação do Serviço Social brasileiro e são de fundamental importância no debate acadêmico e profissional O frágil desempenho dos estudantes nas questões discursivas sugere que possuem um entendimento genérico e, muitas vezes, superficial de conteúdos basilares da formação profissional como definido nas Diretrizes Curriculares Nacionais do Curso de Graduação, tais como : a herança colonial, escravista e patriarcal na formação social, política, cultural e econômica do Brasil ; os movimentos sociais e lutas contra opressões de classe, de gênero, étnico-raciais, de sexualidade, geracionais, entre outras ; estado, questão social e Serviço Social : financeirização do capital, contrarreformas do estado, políticas sociais e seguridade social ; transformações contemporâneas no mundo do trabalho e seus rebatimentos no trabalho da(o) assistente social.
A correção das respostas dos estudantes deixou patente a sua dificuldade de aprofundamento e desdobramento dos processos históricos e concretos em suas particularidades, bem como uma debilidade no enquadramento analítico das temáticas elencadas, recorrendo a interpretações de senso comum e midiáticas, com muitos prejuízos para a compreensão crítica da realidade social.
Há uma acentuada fragilidade analítica, que demonstra limitações na formação profissional dos estudantes de Serviço Social quanto à solidez teórica e metodológica de forma a apresentar uma narrativa pautada pelo rigor científico.
Disso decorre a necessidade de que as unidades de ensino invistam no aprofundamento desses conteúdos para que os estudantes possam articular o saber acadêmico, o significado social da profissão e o seu exercício, para contribuir com a mobilização de sujeitos individuais e coletivos na perspectiva da emancipação humana, bem como para que consigam identificar demandas oriundas da questão social, com base nas particularidades da formação sócio-histórica da sociedade brasileira, formulando respostas profissionais qualificadas e criativas para as mesmas.
Outro aspecto a ser destacado diz respeito à abordagem crítica e analítica dos processos em curso na sociedade brasileira e no mundo, como importante exercício de formação do profissional eticamente comprometido com o enfrentamento da questão social.
Tais abordagens são essenciais para a ancoragem intelectual e interventiva de assistentes sociais em cenário contemporâneo que exige resistência aos avanços conservadores em muitas dimensões da vida social.
Algumas estratégias que podem ser articuladas por órgãos formadores e organizações profissionais são disciplinas eletivas ; seminários e aulas inaugurais ; temáticas em seminários e eventos da categoria ; grupos de pesquisa e socialização de debates e estudos entre as universidades, através de esforços conjuntos de fortalecimento de suas iniciativas.
Quanto ao projeto profissional, vale sublinhar que uma parcela significativa dos estudantes demonstrou uma apropriação restrita, muitas vezes, limitando -o ao seu aspecto normativo ( Código de Ética e leis ) e/ou a"palavras de \textbf{ordem}", jargões e chavões, revelando a falta de reflexão e aprofundamento com base na literatura profissional.
Desse modo, também há a necessidade de que seja aprofundada a discussão do projeto ético e político do Serviço Social no processo formativo, de modo a fortalecer o perfil profissional compromissado com os direitos, a democracia e a liberdade, com a perspectiva da totalidade e com a unidade teoria-prática para análise da questão social e intervenção nas suas expressões, com base na formação profissional generalista e crítica.
E, ainda, que seja empenhado na eliminação de todas as formas de preconceito, discriminação, exploração e opressão, com respeito à diversidade humana, ao meio ambiente e à igualdade social.
O processo de correção permitiu vislumbrar, de modo qualitativo, a formação profissional em Serviço Social no país, e as deficiências identificadas ilustram a precariedade da formação dos futuros profissionais, cujas dificuldades vão desde o desconhecimento das temáticas tratadas nas questões discursivas até a ausência de capacidade argumentativa, de análise e síntese, sinalizando que o desenvolvimento de habilidades e competências previstas nas diretrizes curriculares \textbf{está} profundamente tépido.
Sendo assim, é premente a necessidade de que as unidades de ensino e as entidades profissionais envidem esforços no sentido de desenvolver nos estudantes a capacidade de apreensão teórica da realidade, a partir de mediações concretas apoiadas em temas centrais das Diretrizes Curriculares.
Grande parte das dificuldades dos estudantes se refere ao que, também vem sendo observado no âmbito da graduação em Serviço Social : pouca leitura, fraca apreensão e escrita deficiente.
Observa -se uma linguagem limitada por parte dos estudantes, que nem sempre se valem de conteúdos teórico-conceituais para elaborar os textos de resposta às questões.
Com isso, e em detrimento de um texto de \textbf{ordem} técnica e/ou acadêmica, são fragilizados os conteúdos expostos que, por vezes, são tratados de forma demasiadamente coloquial. coloquial.
É necessário fazer menção aqui, também, às deficiências dos estudantes em formular um texto dissertativo, do ponto de vista da redação, o que, em si, limita, muitas vezes, a compreensão do conteúdo por eles exposto.
Nessa direção, as unidades de ensino devem dar maior atenção ao respeito às regras e normas mais fundamentais da nossa gramática, o que pode ser reforçado pela inclusão ou articulação de disciplinas de língua portuguesa e/ou construção de textos na grade curricular dos cursos.
Como a linguagem – falada e escrita – é um dos principais instrumentos de trabalho dos assistentes sociais –, ela também deve ser priorizada nos Cursos de Graduação em Serviço Social.
Sendo assim, a formulação e organização do pensamento e da argumentação pela escrita precisa de maior atenção por parte das Instituições de ensino em seus projetos pedagógicos, para que possam promover uma formação profissional mais qualitativa, pois o conjunto de dificuldades evidenciadas nas respostas dos estudantes demonstra o rebaixamento da formação profissional que, notadamente, vem se intensificando no decorrer do processo de avaliação do Enade.
Garantir a presença desses conteúdos no processo de formação dos estudantes é essencial para sustentar as intervenções que os futuros profissionais irão desenvolver nos diferentes e desafiadores espaços de trabalho.
O cenário atual exige do assistente social um perfil crítico, criativo, comprometido e competente.
Logo, é o investimento na qualidade da formação profissional oferecida pelos mais diversos Cursos de Graduação em Serviço Social no país e adensada pelo debate profissional promovido pelas entidades representativas da categoria, em âmbito regional e nacional, que vai determinar se, no decorrer do tempo, as lacunas aqui apontadas serão devidamente enfrentadas.

% matched lemmas: apurar, cada, certo, criminalização, estar, lembrar, localização, ordem, torno
\end{textsample}
